\documentclass[11pt,oneside]{article}
\usepackage[T1]{fontenc}
\usepackage{lmodern}
\usepackage{amsmath,amssymb,amsthm,mathtools}
\usepackage[a4paper,margin=28mm,headheight=15pt]{geometry}
\usepackage{microtype}
\usepackage{booktabs,longtable,array,enumitem}
\usepackage[dvipsnames]{xcolor}
\usepackage{xurl}
\usepackage{hyperref}
\usepackage{fancyhdr}
\usepackage{listings}
\hypersetup{colorlinks=true,linkcolor=MidnightBlue,citecolor=MidnightBlue,urlcolor=MidnightBlue,
 pdftitle={Abstract and continuous mathematics through explicit presentations in iPRA},
 pdfauthor={Research report prepared for Vladimir Reshetnikov}}
\newcommand{\N}{\mathbb N}
\newcommand{\Z}{\mathbb Z}
\newcommand{\Q}{\mathbb Q}
\newcommand{\R}{\mathbb R}
\newcommand{\C}{\mathbb C}
\newcommand{\eps}{\varepsilon}
\newcommand{\ipra}{\ensuremath{\mathsf{iPRA}}}
\newcommand{\PR}{\mathsf{PR}}
\newcommand{\norm}[1]{\left\lVert #1\right\rVert}
\newcommand{\abs}[1]{\left\lvert #1\right\rvert}
\newcommand{\ip}[2]{\left\langle #1,#2\right\rangle}
\newcommand{\D}{\mathcal D}
\newcommand{\Reg}{\operatorname{Reg}}
\newcommand{\Eq}{\operatorname{Eq}}
\newcommand{\Net}{\operatorname{Net}}
\newcommand{\Lip}{\operatorname{Lip}}
\newcommand{\ceil}[1]{\left\lceil #1\right\rceil}
\newcommand{\floor}[1]{\left\lfloor #1\right\rfloor}
\newcommand{\repo}{\texttt{ProveIt}}
\newcommand{\sourcepath}[1]{\texttt{\small\path{#1}}}
\newtheorem{theorem}{Theorem}[section]
\newtheorem{lemma}[theorem]{Lemma}
\newtheorem{proposition}[theorem]{Proposition}
\newtheorem{corollary}[theorem]{Corollary}
\theoremstyle{definition}
\newtheorem{definition}[theorem]{Definition}
\newtheorem{example}[theorem]{Example}
\theoremstyle{remark}
\newtheorem{remark}[theorem]{Remark}
\newtheorem{warning}[theorem]{Boundary}
\setlist{itemsep=3pt,topsep=6pt}
\setlength{\emergencystretch}{3em}
\pagestyle{fancy}
\fancyhf{}
\fancyhead[L]{\small Explicit presentations in \(\ipra\)}
\fancyhead[R]{\small Research report}
\fancyfoot[C]{\thepage}
\lstset{basicstyle=\ttfamily\footnotesize,breaklines=true,columns=fullflexible,
 frame=single,rulecolor=\color{black!20},backgroundcolor=\color{black!2},
 showstringspaces=false}
\begin{document}
\hypersetup{pageanchor=false}
\begin{titlepage}
\vspace*{13mm}
{\color{MidnightBlue}\large FOUNDATIONS AND CONSTRUCTIVE ANALYSIS\par}
\vspace{10mm}
{\raggedright\huge\bfseries Abstract and continuous mathematics through explicit presentations in \(\ipra\)\par}
\vspace{8mm}
{\Large Completion, function spaces, continuous logic, coercive equations, and analytic approximation\par}
\vspace{13mm}
{\large Prepared for Vladimir Reshetnikov\\20 November 2026\par}
\vspace{10mm}
\noindent\textbf{Main contribution.}
A common interpretation discipline is developed into concrete continuous mathematics:
completion of presented metric spaces; quantitative function-space approximation and
gluing; finite-net evaluation of continuous formulas; a constructive coercive-operator
solver with an explicit clock; a variable-coefficient weak Dirichlet problem; and
simultaneous approximation of complex derivatives. The report also proves precise
limits on what the available data can determine.

\vspace{7mm}
\noindent\textbf{Verification boundary.}
The mathematical arguments are human proofs with explicit arithmetic reductions.
The accompanying small Lean development checks selected structural and quantitative
lemmas. The new main analytic theorems have \emph{not} yet been compiled into
\(\texttt{Prf PRAAx}\). Complete theorem targets, required definitions, and an ordered
formalization plan are provided.

\vfill
\noindent Repository review pinned to\\
\texttt{5f3be02388442f16487d2dc5c2529961dd783d2b}.\\[2mm]
The repository was read, not modified.
\end{titlepage}
\hypersetup{pageanchor=true}

\begin{abstract}
First-order intuitionistic primitive recursive arithmetic cannot simply be given
variables ranging over arbitrary real functions, all metric spaces, or all models.
Nevertheless, a substantial part of continuous mathematics admits a precise
arithmetical treatment once a presentation specifies finite approximants, uniform
moduli, and quantitative stability. We organize this treatment into three levels:
finite arithmetic, coherent represented objects, and semantic interpretation.
The arithmetic level is the target of actual \(\ipra\) derivations; the other two
levels explain which classical or constructive theorem those derivations express.
We apply the discipline to categories of completions, function spaces, metric model
theory, an infinite-dimensional coercive equation and a one-dimensional elliptic
boundary-value problem, and families of holomorphic functions with derivative
control. We separate effective total boundedness from subsequence selection,
strong Hilbert names from weak functional data, and analytic bounds from a merely
qualitative analyticity assertion. Every positive result has an explicit algorithm
and a finite error estimate. Every negative conclusion is stated at its actual
metatheoretic strength.
\end{abstract}
\tableofcontents
\newpage
\section{Scope, baseline, and the interpretation discipline}
\label{sec:foundations}

\subsection{What the existing project already supplies}

The review is pinned to commit
\texttt{5f3be02388442f16487d2dc5c2529961dd783d2b}, rather than to a moving
branch. The core README distinguishes research packets from checked proof
libraries. In particular, a packet's statement that its proof is intended for
\(\ipra\) is not itself a checked derivation. The existing source goes
considerably beyond the early packet descriptions: it already contains
presented-sort interpretation, observation composition, finite nets, coherent
minimizers, and many finite analytic estimates. We reuse this architecture
instead of proposing another renaming of a certificate interface
\cite{repoMain,repoSetoid,repoExtensional,repoFramework,repoCoherent}.

The most relevant existing mechanisms are as follows. Paths in this report are
relative to \sourcepath{Logic/PrimitiveRecursiveArithmetic/}.
\begin{center}
\small
\begin{tabular}{p{.27\textwidth}p{.64\textwidth}}
\toprule
Existing mechanism & Consequence for the present work\\
\midrule
Setoid interpretation &
Source equality may be interpreted by an equivalence formula; it need not be
a decidable equality test. Quantifiers are relativized to a domain formula.\\
Precision composition &
The demand of a composite is already derived by composing demands. This
logical rule does not manufacture a missing modulus.\\
Coherence and stability &
Independent approximate witnesses require a stability or selection argument
before they define one continuous object.\\
Finite nets and covers &
A finite net suffices for approximate extrema and robust observations.
An arbitrary open cover or an infinite sequence has additional requirements.\\
Analytic charts and contour packets &
Radius, magnitude, separation, differentiability and tail data are explicit.
They cannot be erased merely by calling an input analytic.\\
\bottomrule
\end{tabular}
\end{center}

The new work develops \emph{instances that exercise these distinctions}:
function spaces rather than only individual functions; quantification over
metric sorts rather than only point evaluation; coercive equations in an
infinite-dimensional completion; and simultaneous control of derivatives.
The individual mathematical ideas are classical or constructive-analysis
tools. No historical priority claim is made. The contribution is their
explicit, mutually compatible arithmetic formulation, the detailed proofs,
and the diagnosis of missing information.

\subsection{The exact base theory}

\begin{definition}[The arithmetic theory]
\(\ipra\) has one numerical sort, intuitionistic first-order logic with equality,
a symbol for each specified primitive-recursive definition, its defining
equations, the successor axioms, and induction for quantifier-free formulas.
Free numerical parameters are allowed in induction.
\end{definition}

A primitive-recursive definition is built from zero, successor, projections,
composition, and primitive recursion. A finite tuple, list, polynomial or
matrix is encoded by a natural number. The code predicates and all component
operations used below are primitive recursive. Integers are signed natural
codes; rational numbers are pairs consisting of an integer numerator and a
positive natural denominator. A malformed code is assigned a harmless default
value, and validity is kept as a separate premise. Equality and order of
rationals are decided by multiplication by positive denominators. Reducing
fractions to lowest terms is optional.

A statement involving \(2^{-k}\), a rational power, a finite maximum, or an
integer ceiling therefore abbreviates a finite arithmetic formula. Bounded
search means checking every candidate below an explicit numerical bound, not
searching until an unspecified terminating event occurs. No efficiency claim
is implicit: bounds such as \(2^{2^k}\) are primitive recursive.

\begin{lemma}[Bounded conjunctions require only quantifier-free induction]
\label{lem:boolean}
Let \(P(i,\vec a)\) be a decidable primitive-recursive predicate. Its bounded
universal and existential assertions can be represented by primitive-recursive
Boolean folds. A loop whose invariant consists of finitely many such assertions
can be verified by quantifier-free induction.
\end{lemma}
\begin{proof}
Choose a Boolean characteristic \(p(i,\vec a)\in\{0,1\}\). Define
\(b(0,\vec a)=1\) and
\(b(n+1,\vec a)=b(n,\vec a)p(n,\vec a)\).
The finite scan has a quantifier-free state invariant recording the accumulated
Boolean and any requested component index. Its defining equations give the
bounded introduction and elimination rules: a product of Boolean values is
\(1\) exactly when every factor is \(1\). To prove the backward implication
one inducts on the scan while the required bounded universal premises are
available in the context, instantiating the next premise at each step.
To prove the forward implication fix the requested index first and use
quantifier-free induction to preserve its implication from the accumulator.
Existential search is the analogous disjunction fold with a stored least
witness. Bounds and finite tuples can be included in the state.
This does not allow an unbounded universal formula to become an induction
invariant: each unbounded query is first fixed as a parameter.
\end{proof}

\begin{lemma}[Rational separation]
\label{lem:separation}
For rationals \(r,s,C\) with \(C\geq0\), if
\(r\leq s+C2^{-k}\) for every \(k\), then \(r\leq s\).
The conclusion has a direct intuitionistic proof from the displayed family
of finite inequalities.
\end{lemma}
\begin{proof}
Rational order is decidable. If \(r>s\), write the positive difference as
\(a/b\), with integers \(a,b\geq1\). Choose an integer \(L\geq Cb+1\), and
put \(k=L+1\). The elementary induction \(2^L\geq L+1\) gives
\(C2^{-k}<1/b\leq a/b\), contradicting the \(k\)-th inequality. No real
limit rule or Markov principle is being used.
\end{proof}

\subsection{Three levels that must not be conflated}

\paragraph{Finite arithmetic.}
A theorem at this level is a sentence or open sequent in the actual arithmetic
language. Examples are a squared rational norm estimate, a bound on a finite
iteration, or the correctness of a bounded scan. This is the layer whose
proof terms must eventually have type
\[
                  \Gamma\vdash[\mathsf{PRAAx}]\varphi.
\]

\paragraph{Coherent presentations.}
An infinite object is described by an externally specified primitive-recursive
family of approximants. Its regularity and equality are \emph{formulas}, usually
with unbounded quantifiers, whose individual instances reduce to finite
arithmetic. Operations come with finite demand and error estimates.

\paragraph{Semantic interpretation.}
One may explain a regular rational name as a real number, a regular sequence
of rational polynomials as a continuous function, or an energy Cauchy name
as a weak solution. This explanation is carried out in an ambient constructive
metatheory. It does not add real-number variables or a completeness axiom to
first-order \(\ipra\).

\begin{definition}[A regular metric name]
\label{def:regular}
Let \(X_0\) be a coded countable set with a rational-valued metric \(d\), or
with rational squared norms sufficient to express the metric inequalities.
A regular name is a specified map \(p:\N\to X_0\) satisfying
\[
       d(p(n),p(m))\leq2^{-n}\qquad(m\geq n).
\]
Two names represent the same point when their distance tends to zero with
the usual rational precision tests. For regular names one convenient test is
\[
       d(p(n+2),q(n+2))\leq2^{-n}\qquad\text{for all }n.
\]
\end{definition}

The last test agrees with equality in a metric completion. Its transitivity
uses finer indices and rational separation, not a false assertion that
two errors of size \(2^{-n}\) still sum to \(2^{-n}\).
For example, if \(p\sim q\sim r\), then for a fixed index \(n\), arbitrarily
fine indices \(j\geq n+2\) give
\[
d(p(n+2),r(n+2))
 \leq 2^{-n-2}+2^{3-j}+2^{-n-2}.
\]
Separation gives a bound \(2^{-n-1}\), hence the required \(2^{-n}\).
The same argument proves independence of bounded shifts used to normalize
names. When only squared distances are rational, compare their rational
squares, or first use rational upper bounds with an arbitrary error.

The space of primitive-recursive names is not claimed complete under
arbitrary externally given Cauchy sequences. It is closed under
\emph{jointly primitive-recursive} Cauchy constructions with a supplied
primitive-recursive modulus: the diagonal name uses finitely many fixed
operations and the modulus. That is exactly the closure used below.

\subsection{Schemas, numerical parameters, and the universal-evaluator trap}

A fixed producer \(p(a,n)\) may depend on a natural parameter \(a\). It gives
a numerical presented sort with domain
\(\Reg_p(a)\) and an equality formula \(\Eq_p(a,b)\). The existing
\texttt{SetoidInterpretation.sound} can transport a first-order theory
through such a presentation once its domain, congruence and axiom obligations
are actually proved.

There is also a broader external statement: \emph{for every chosen
primitive-recursive producer and every chosen list of proved contracts, the
following construction works}. This is a theorem schema, not a single
first-order quantifier over programs. A finite library of operation symbols
can have a primitive-recursive syntax evaluator when its recursion is
appropriately bounded; this does not give a universal evaluator for all
primitive-recursive functions.

The obstruction is elementary. If a total primitive-recursive
\(U(e,n)\) evaluated every unary primitive-recursive function, then
\(g(n)=U(n,n)+1\) would itself have an index \(e\), giving
\(g(e)=U(e,e)+1=g(e)+1\). Thus one must not infer that a countable collection
of names can be packaged into one such evaluator. The repository already
records this issue \cite{repoUniversal}. In this report the categories and
function-space constructions are external categories of specified
presentations, with object-theory instances for each fixed signature.

\subsection{A restricted compilation theorem}

\begin{theorem}[Compilation of the finite parts]
\label{thm:compilation}
Fix the primitive-recursive symbols used by one of the constructions below.
Suppose its input contracts are supplied as \(\ipra\) derivations. Suppose
also that every loop has a specified quantifier-free or Booleanized bounded
invariant, and that the displayed arithmetic inequalities proving its base
and step have been derived. Then the finite correctness statement of the
construction has an \(\ipra\) derivation. The transformation is effective
on these supplied proof terms.
\end{theorem}
\begin{proof}
Composition substitutes proved equations and uses ordinary intuitionistic
cut. Primitive recursion uses its two defining equations. A loop adds the
quantifier-free induction instance for its invariant, substitutes the proved
base and step, and specializes the result at the explicit execution length.
Bounded finite checks use Lemma~\ref{lem:boolean}. Rational inequalities are
cleared with positive denominators before this process. An arbitrary
observation index is treated as a free parameter throughout and universally
generalized only after the finite derivation is complete. A finite collection
of such compilations is composed by cut. All of these are rules already
available in the object calculus.
\end{proof}

This theorem is deliberately restricted. It does not assert that every true
primitive-recursive specification has a proof, that checking a proof yields
its truth inside the same theory, or that every constructive analytic proof
has primitive-recursive witnesses. In particular, a host Lean proof about
\(\N\) or \(\Q\) is not automatically such an object derivation.

\subsection{What counts as resolving a difficulty}

The report uses three distinct forms of progress.
\begin{enumerate}
\item \textbf{Producing missing data.} A finite net produces a positive
denominator for interpolation; coercivity produces a contraction gap and
an iteration clock; a radius gap produces a tail modulus for derivatives.
\item \textbf{Proving a quantitative replacement.} Uniform total boundedness
and approximate quantifier evaluation replace a stronger qualitative
compactness assertion at the observation level actually needed.
\item \textbf{Identifying an irreducible input.} A Riesz tail modulus or
analytic radius cannot generally be recovered from weaker supplied data.
An impossibility proof prevents an unsound promise of a universal repair.
\end{enumerate}
The distinctions are mathematically substantive. Merely assuming the desired
conclusion as a certificate would not resolve the corresponding problem.

\section{Completion as a categorical construction with arithmetic interfaces}
\label{sec:cm-completion}

The existing repository already has a substantial abstract interface.  In
\path{Foundations/SetoidInterpretation.lean}, numerical domain and equality
formulas interpret a source theory, and the theorem \texttt{sound} transports
its intuitionistic derivations once closure, congruence, and interpreted axioms
have been proved.  \texttt{PresentedCategory.lean} and \texttt{Yoneda.lean}
supply actual arithmetic derivations for specified categorical presentations.
Consequently the missing step is usually not the equational manipulation of a
category.  It is constructing a suitable presentation of its continuous
objects and proving that its proposed operations preserve those objects.
This section supplies one such step, without asserting that every analytic
category has a numerical universe.

\subsection{Two levels that must remain distinct}

A \emph{rational metric presentation} consists of a decidable set of numerical
codes $D$, a distinguished valid code, and a fixed primitive recursive rational
pseudometric $d:D\times D\to\Q_{\geq0}$, with proofs of its pseudometric laws.
The use of a pseudometric permits different rational codes to denote the same
point.  Rational arithmetic means coded integer numerators and positive
natural denominators; all inequalities below can be cleared of denominators.
There is no assumption that $D$ is finite or discrete in its metric.

Externally, its completion $\widehat D$ has the usual mathematical meaning.
Internally, a family of names is given by a \emph{specified} PR function
$q(p,n)\in D$, where $p$ is a finite tuple of numerical parameters.  Its
validity is an arithmetic formula in $p$.  A different producer $q'$ is a
different instance of the scheme, unless a fixed uniform PR family containing
both has actually been supplied.  We do not form a numerical sort of all PR
programs equipped with a total PR evaluator.  The repository's
\texttt{UniversalEvaluatorObstruction.lean} is directly relevant to this
restriction.

\begin{definition}[Fast names and equality]
\label{def:cm-fast}
A fast name is a sequence $q_n\in D$ satisfying
\[
             n\leq m\quad\Longrightarrow\quad d(q_n,q_m)\leq 2^{-n}.
\tag{C1}\label{eq:cm-fast}
\]
Two fast names are equivalent if
\[
 q\sim r\quad\Longleftrightarrow\quad
 \forall k\ \exists N\ \forall n\geq N\quad d(q_n,r_n)\leq 2^{-k}.
\tag{C2}\label{eq:cm-eq}
\]
For fixed PR producers these are first-order arithmetic formulas, not
quantification over functions.
\end{definition}

Reflexivity and symmetry of $\sim$ follow from the metric laws.  For
transitivity, instantiate the two equivalences at $k+1$, take the maximum of
the resulting two indices, and apply the triangle inequality.  This uses
intuitionistic existential elimination and introduction and contains no choice
principle.  Constant names embed $D$.  The rational sequence
$d(q_{n+1},r_{n+1})$ is fast, since the reverse triangle inequality gives, for
$m\geq n$,
\[
 |d(q_{n+1},r_{n+1})-d(q_{m+1},r_{m+1})|
 \leq 2^{-n-1}+2^{-n-1}=2^{-n}.
\]
It therefore supplies the completed distance in the semantic interpretation.
A fast name is at distance at most $2^{-n}$ from its $n$th constant
approximation.  These observations explain the completion notation.  They do
not imply that an arbitrary sequence of PR names has a PR limit; a uniform
producer and an effective outer Cauchy rate are needed for that assertion.

\subsection{A dense evaluator extends with one precision substitution}

Let $D_X,D_Y$ be two presentations.  A dense evaluator is a fixed PR function
$E:D_X\times\N\to D_Y$.  Its output need only approximate the value at a dense
input; square roots, special functions, and bounded linear operators can be
handled this way.  Let $\mu:\N\to\N$ be a specified PR function, nondecreasing
and satisfying $\mu(k)\geq k$.  The essential hypothesis is the following
finite, rational inequality:
\[
\begin{split}
 d_X(a,b)\leq 2^{-\mu(k)},\quad r,s\geq k+2
 \quad\Longrightarrow\quad
 d_Y(E(a,r),E(b,s))\leq2^{-k}.
\end{split}
\tag{C3}\label{eq:cm-joint}
\]
The theorem uses a proof of this implication; it does not test its universal
validity.  When $E(a,r)$ is known to approximate a map $f(a)$ within $2^{-r}$
and $f$ has spatial modulus $\mu_0$, the normalized choice
\[
     \mu(k)=\max\{k,\mu_0(0),\ldots,\mu_0(k+1)\}
\]
suffices: the three contributions are $2^{-k-2}$, $2^{-k-1}$,
and $2^{-k-2}$.  Thus \eqref{eq:cm-joint} is a convenient way to combine
spatial continuity and evaluation error.

\begin{theorem}[Dense extension with explicit demand]
\label{thm:cm-extension}
Under \eqref{eq:cm-joint}, define
\[
             (\widehat E q)_n=E(q_{\mu(n)},n+2).
\tag{C4}\label{eq:cm-extension}
\]
This is a fast name.  Equivalent input names give equivalent output names.
For a specified PR input family it is a specified PR output family, and both
assertions are first-order schemes over the two families.  Semantically this
construction is a continuous extension to the completions, and it is the
unique uniformly continuous extension of its values on the dense points.
\end{theorem}

\begin{proof}
If $m\geq n$, monotonicity gives $\mu(m)\geq\mu(n)$, so
\eqref{eq:cm-fast} gives
$d_X(q_{\mu(n)},q_{\mu(m)})\leq2^{-\mu(n)}$.
The output precisions $n+2,m+2$ satisfy the two guards in
\eqref{eq:cm-joint}.  Its instance at $k=n$ proves
\[
      d_Y((\widehat E q)_n,(\widehat E q)_m)\leq2^{-n}.
\]
This is the whole fast-name calculation.

Now suppose $q\sim r$, and fix an output challenge $k$.  Use
\eqref{eq:cm-eq} at input precision $\mu(k)$, obtaining $N$.  For
$n\geq\max(N,k)$, the inequality $\mu(n)\geq n$ makes
$d_X(q_{\mu(n)},r_{\mu(n)})\leq2^{-\mu(k)}$.  Both evaluator precisions
are $n+2\geq k+2$.  Equation \eqref{eq:cm-joint} at the fixed challenge $k$
therefore gives the required output distance $\leq2^{-k}$.  This proves
extensionality without selecting a function of the existential witnesses in
\eqref{eq:cm-eq}.

For continuity in the semantic completion, give the input distance a spare
factor of two.  If
$d_{\widehat X}(x,y)\leq2^{-\mu(k+1)-1}$, sufficiently late approximations
have distance at most $2^{-\mu(k+1)}$.  Applying
\eqref{eq:cm-joint} to those approximations and passing to the named limit
gives output distance at most $2^{-k-1}$, hence at most $2^{-k}$.  Thus
$k\mapsto\mu(k+1)+1$ is a safe modulus for the extension.  This semantic
sentence can equally be replaced by its finite rational inequalities.

For uniqueness, suppose $F,G:\widehat X\to\widehat Y$ are uniformly
continuous and agree on constant names.  Given $x=[q]$ and $k$, choose a dense
approximation $q_N$ close enough to $x$ for both output errors to be at most
$2^{-k-1}$.  Then
\[
 d(Fx,Gx)\leq d(Fx,Fq_N)+d(Fq_N,Gq_N)+d(Gq_N,Gx)\leq2^{-k}.
\]
The middle term is zero.  Since this holds at every precision, the two output
names are equivalent.  For specified moduli, $N$ is the maximum of their
explicit input demands; no compactness argument is hidden here.
\end{proof}

For exact dense maps $f:D_X\to D_Y$ with modulus $\mu_f$, the simpler lift
$(\widehat f q)_n=f(q_{\mu_f(n)})$ is already fast.  The identity has
$\mu_{\mathrm{id}}(n)=n$.  For $g\circ f$, use
$\mu_{g\circ f}=\mu_f\circ\mu_g$; the two composite lifted expressions are
literally identical.  Approximate dense evaluators compose by substitution of
the corresponding name transformations, and uniqueness identifies different
valid schedules for the same dense map.  In this precise external or
schematic sense completion is functorial.  It is not an internal category
whose objects are all possible PR implementations.

\begin{proposition}[Products and change of dense presentation]
\label{prop:cm-products-density}
The completion construction preserves finite products equipped with the
maximum metric.  Moreover, suppose $f:D_X\to D_Y$ is a specified PR isometry
and a PR selector $s:D_Y\times\N\to D_X$ satisfies
\[
                    d_Y(f(s(b,k)),b)\leq2^{-k}.
\tag{C5}\label{eq:cm-dense-selector}
\]
Then $\widehat f$ is an isomorphism of the presented completions, with an
explicit inverse on names.
\end{proposition}

\begin{proof}
A name in the maximum product metric is fast exactly when both coordinate
names are fast.  Equivalence in the product is equivalent to equivalence in
both coordinates: the forward implication follows by projection; for the
reverse implication take the maximum of the two eventual indices.  Pairing
and projection therefore give inverse name transformations.  The terminal
one-point presentation covers the empty product.

For the second assertion, define
\[
                  (Hy)_n=s(y_{n+2},n+2).
\tag{C6}\label{eq:cm-density-inverse}
\]
For $m\geq n$, the isometry law and three applications of the triangle
inequality give
\[
\begin{split}
 d_X((Hy)_n,(Hy)_m)
 &=d_Y(f(Hy)_n,f(Hy)_m)\\
 &\leq2^{-n-2}+d_Y(y_{n+2},y_{m+2})+2^{-m-2}\\
 &\leq3\,2^{-n-2}\leq2^{-n}.
\end{split}
\]
Here $f(Hy)_n$ denotes $f((Hy)_n)$.  Thus $H$ produces a fast name.  Also
$d_Y(f((Hy)_n),y_{n+2})\leq2^{-n-2}$, so $\widehat f(Hy)\sim y$.
For an $X$-name $x$, the isometry law similarly gives
$d_X((H\widehat f x)_n,x_{n+2})\leq2^{-n-2}$, hence
$H\widehat f x\sim x$.  Finally, if $y\sim z$, applying the triangle
inequality to $f((Hy)_n),y_{n+2},z_{n+2},f((Hz)_n)$ proves $Hy\sim Hz$.
All output indices in these estimates are displayed; for example an
$n$-stage error bounded by $5\,2^{-n-2}$ is below $2^{-k}$ once
$n\geq k+1$.  Thus the two inverse equations are first-order equality
statements about the supplied name families.
\end{proof}

Condition \eqref{eq:cm-dense-selector} is positive information.  An unadorned
assertion that $f(D_X)$ is dense need not provide a PR approximation selector.
This is the continuous analogue of the finite-cover versus cardinality-bound
distinction already identified in the repository's categorical packet.
The proposition permits safe transport between, for example, rational and
dyadic descriptions of a compact interval, or two explicit dense descriptions
of a normed space, once the displayed selector is proved.

\subsection{A direct Lean implementation boundary}

The first implementation should keep codes, rational distances, and producers
explicit.  One can use the following metatheoretic records and then replace
the fields asserting inequalities by the corresponding derivations in
\texttt{Prf PRAAx}:
\begin{verbatim}
structure MetricPresentation where
  valid : Nat -> Bool
  dist  : Nat -> Nat -> RatCode
  -- Proofs: closure, nonnegativity, symmetry, triangle.

def Fast (D : MetricPresentation) (q : Nat -> Nat) : Prop :=
  (forall n, D.valid (q n) = true) /\
  (forall n m, n <= m -> D.dist (q n) (q m) <= dyadic n)

structure DenseEvaluator (X Y : MetricPresentation) where
  eval : Nat -> Nat -> Nat
  mu   : Nat -> Nat
  -- PR realizers and monotonicity proofs are separate fields.
  joint : forall a b k r s, ValidInputs a b ->
    X.dist a b <= dyadic (mu k) -> k + 2 <= r -> k + 2 <= s ->
    Y.dist (eval a r) (eval b s) <= dyadic k

def extend (E : DenseEvaluator X Y) (q : Nat -> Nat) (n : Nat) :=
  E.eval (q (E.mu n)) (n + 2)
\end{verbatim}

This block is a blueprint, not compilable source and not an assertion that
\texttt{Nat -> Nat} is an object-language sort.  In the arithmetic layer it is
replaced by an explicit PR term $q(p,n)$ and by formulas
$\mathrm{Fast}_q(p)$ and $\mathrm{Eq}_{q,r}(p,p')$.  The principal target is
\[
 \ipra\vdash\forall p\bigl(\mathrm{Fast}_q(p)
                  \to\mathrm{Fast}_{\widehat E q}(p)\bigr),
\]
with the modulus and joint-law hypotheses supplied as derivations.  Its proof
instantiates those hypotheses at arbitrary terms and combines rational
inequalities; it does not induct on $\mathrm{Fast}_q$, which contains
unbounded quantifiers.  Product and density-selector targets use the same
procedure.  The existing setoid-interpretation theorem then applies to a
\emph{fixed} numerical family closed under the operations that have actually
been defined.

\section{A finite-observation semantics for compact metric structures}
\label{sec:cm-model}

Continuous model theory uses metric structures, uniformly continuous symbols,
real-valued predicates, and $\inf/\sup$ quantifiers; this is the setting
surveyed by Ben Yaacov, Berenstein, Henson, and Usvyatsov
\cite{bbhu2008}.  The following development supplies a deliberately narrower
arithmetic realization with explicit total-boundedness data.  It preserves
quantifier order, including alternation.  It neither imports the general
compactness theorem of continuous logic nor replaces it by a finite-model
claim.

\subsection{Language, moduli, and finite restrictions}

Fix finitely many metric sorts and a finite signature. In the exact dense
version, constants have supplied dense codes; named constants are covered
by the approximate-evaluation variant below. For simplicity first suppose
that function symbols act by fixed
PR maps between the dense presentations and that the predicates are fixed
PR rational functions taking values in $[0,1]$.  Each symbol is accompanied
by a PR modulus and proofs of its continuity inequalities.  The dense
extension theorem interprets the symbols on the completions.

Terms are variables, constants, and finite applications of function symbols.
An atomic expression is a predicate applied to terms, including the bounded
metric predicate.  Formulas are built from atoms and rational constants in
$[0,1]$ by $\min$, $\max$, $1-u$, and the quantifiers
$\inf_{x:s}$ and $\sup_{x:s}$.  With the convention that $0$ is true,
$\max$ plays the role of conjunction and $\min$ that of disjunction.
The operation $1-u$ is an arithmetic connective in the object being studied;
its use does not add classical excluded middle to the surrounding iPRA
proof.

A modulus of a term or atom is computed by composition.  For a variable use
$\alpha(k)=k$.  If $f$ has modulus $\mu_f$ and the arguments have moduli
$\alpha_i$, use
\[
            \alpha_{f(t_1,\ldots,t_r)}(k)
                    =\max_i\alpha_i(\mu_f(k)).
\tag{M1}\label{eq:cm-atom-modulus}
\]
Use the same rule for a predicate.  The metric has safe modulus $k+1$ with
respect to the maximum distance on its two arguments.  Constants have no
input demand.  Extend the modulus trivially to unused variables.  For a fixed
formula $\varphi$, let $\alpha_\varphi(k)$ be the maximum of the moduli of
all its atomic occurrences, normalized to be nondecreasing and at least $k$.
Then $\delta=2^{-\alpha_\varphi(k)}$-close assignments make every atomic
value differ by at most $2^{-k}$.

For each sort choose a finite nonempty list $A_s$ of dense points.  The finite
restriction $V_A(\varphi,\bar a)\in\Q\cap[0,1]$ is defined recursively:
use the dense evaluations at atoms, the indicated rational operations at
connectives, and the finite minimum or maximum over $A_s$ at a quantifier of
sort $s$.  Bound variables have their usual scope; the recursion does not
interchange two quantifiers.  Terms continue to be evaluated in the ambient
dense presentation.  Thus $A_s$ need not be closed under a function symbol.
In particular $V_A$ is not being advertised as a finite submodel of the
original structure.

\begin{theorem}[Quantitative back-and-forth for finite restrictions]
\label{thm:cm-backforth}
Let $A_s,B_s$ be finite nonempty lists for every sort.  Suppose each point of
$A_s$ is within $\delta$ of some point of $B_s$, and conversely.  Suppose
also that corresponding free-variable assignments are coordinatewise within
$\delta$, and that every atomic occurrence of $\varphi$ changes by at most
$\eta$ on $\delta$-close assignments.  Then
\[
                |V_A(\varphi,\bar a)-V_B(\varphi,\bar b)|\leq\eta.
\tag{M2}\label{eq:cm-transfer}
\]
The error does not acquire a factor from the number of quantifier alternations.
If atomic values are themselves evaluated with errors at most $e_A,e_B$, the
right side becomes $\eta+e_A+e_B$.
\end{theorem}

\begin{proof}
We use structural induction on the fixed formula to produce its arithmetic
proof.  The atom case is the hypothesis.  For $1-u$, the error is unchanged.
For $\min$ and $\max$, the elementary inequality
\[
 |\max(u_1,u_2)-\max(v_1,v_2)|
 \leq\max\{|u_1-v_1|,|u_2-v_2|\}
\]
and its minimum counterpart preserve the same bound.  Each follows by the
two inequalities $u_i\leq v_i+\eta$ and $v_i\leq u_i+\eta$; taking the
appropriate maximum or minimum proves both directions.

Consider $\inf_{x:s}\psi$.  For any $a\in A_s$, choose a $b\in B_s$
within $\delta$.  The extended assignments still satisfy the same
coordinatewise relation.  The induction hypothesis gives
\[
 V_A(\psi,\bar a,a)\geq V_B(\psi,\bar b,b)-\eta
             \geq\min_{b'\in B_s}V_B(\psi,\bar b,b')-\eta.
\]
Taking the finite minimum over $a$ proves one inequality.  The reverse
covering direction proves the other, with $A$ and $B$ exchanged.  Alternatively,
a first minimizing entry can be selected by a bounded scan and used as the
witness.  For $\sup$, replace minima by maxima and reverse the inequalities.
The error at a child is not added to a new rounding error: the corresponding
bound-variable choices were already included in its atom-level hypothesis.
This explains the absence of a depth factor.

Finally, evaluating every atom within an error $e$ changes the entire finite
formula by at most $e$, by the identical induction with the same domains.
Applying this on both sides of \eqref{eq:cm-transfer} gives
$e_A+\eta+e_B$.  Only rational operations and finite minima, maxima, and
searches appear in this proof.
\end{proof}

The finite covering relation is a bounded formula: after fixing the lists,
all its existential witnesses have an explicit list-length bound.  It may be
Booleanized.  Induction used to compute a minimum is on a finite prefix with
its rational lower-bound invariant.  Thus the proof does not require induction
on a formula involving quantification over an entire completion.  An external
induction on the fixed source syntax generates a finite iPRA proof, much as
in the existing interpretation theorem.

\subsection{From nets to names of formula values}

Assume each sort has a fixed PR net producer $N_s(t)$, of positive PR length,
with the supplied covering theorem
\[
       \forall a\in D_s\ \exists j<|N_s(t)|\quad
                   d_s(a,N_s(t)_j)\leq2^{-t}.
\tag{M3}\label{eq:cm-nets}
\]
The unbounded quantifier over $a$ in this theorem is a hypothesis already
proved; validity of a purported net producer is not decided by inspecting one
finite list.  This is a quantitative total-boundedness assumption, stronger
than merely calling a sort bounded.

\begin{theorem}[PR values for compactly presented metric formulas]
\label{thm:cm-values}
Fix the signature, its PR producers, and a formula $\varphi$.  For specified
fast input names $q^1,\ldots,q^r$, put
\[
 v_n=V_{N(\alpha_\varphi(n))}
       \bigl(\varphi,q^1_{\alpha_\varphi(n)},\ldots,
                          q^r_{\alpha_\varphi(n)}\bigr).
\tag{M4}\label{eq:cm-formula-name}
\]
Then $v_n$ is a PR fast Cauchy name.  Its represented value is independent of
valid input names and of the chosen net producers.  In the semantic
completion it gives the usual continuous $\inf/\sup$ interpretation of the
formula.  Each instance has a first-order iPRA realization by finite
inequalities.
\end{theorem}

\begin{proof}
For $m\geq n$, the free-variable tuples in \eqref{eq:cm-formula-name} are
within $\delta=2^{-\alpha_\varphi(n)}$ by fastness and monotonicity.  Every
point of the coarser net is a dense point, so the finer net covers it within
$2^{-\alpha_\varphi(m)}\leq\delta$; the coarser net covers every point of
the finer net within $\delta$ as well.  Theorem~\ref{thm:cm-backforth} gives
$|v_n-v_m|\leq2^{-n}$.  This is an explicit Cauchy proof; an independently
assumed existence of a real infimum is unnecessary.

The same comparison proves independence from two choices of net producers.
For equivalent input names, fix an output precision, instantiate the
input equivalences at the required atomic precision, and take stages beyond
the maximum of the resulting finitely many indices.  Applying
Theorem~\ref{thm:cm-backforth} at that fixed precision shows that the two
output names are equivalent.

To identify the quantifier semantics, it is enough to handle an infimum.
Fix the other variables.  For any dense $a$, the net theorem supplies points
within the required atomic modulus of $a$; their formula values bound the
finite minimum from above up to the displayed error.  Hence the named value
is a lower bound for all dense-input values.  Conversely every finite
minimum is attained at a net entry.  Those entries give values arbitrarily
close above the named lower bound.  Continuity extends the lower-bound
inequality from the dense points to arbitrary named points.  This is exactly
the defining lower-bound and approximate-attainment property of the infimum.
The supremum argument is dual.  Notice that the approximate minimizing net
entry may change with precision.  No coherent convergent sequence of
minimizers has been asserted.

Every evaluator in \eqref{eq:cm-formula-name} uses fixed PR symbols and a
finite syntax tree.  Evaluation at a quantifier is a bounded scan through a
PR-length list; nested scans have a fixed finite nesting depth.  Consequently
the displayed producer is PR.  The arithmetic proof applies the covering
theorems only to the particular finite entries used by the scans, and uses
bounded invariants for those scans.  It does not induce on a real-valued
satisfaction predicate.
\end{proof}

A useful version permits PR rational approximations to dense atomic values.
Evaluate all atoms to error $2^{-n-2}$, use nets and free-variable stages at
$\alpha_\varphi(n+2)$, and retain the same finite evaluation.  Comparing
stages $n\leq m$ gives an error at most
\[
                2^{-n-2}+2^{-n-2}+2^{-m-2}\leq 2^{-n}.
\]
Thus the conclusion includes predicates defined by analytic evaluation
algorithms, not only rational polynomials.  For an approximate dense function
symbol, first compose its evaluator with the finitely many surrounding term
and predicate evaluators, allocating precision along that fixed term tree.
The resulting atomic evaluator and atomic modulus are the explicit inputs to
the theorem.

The theorem gives a practical model-theoretic transfer principle.  If two
compact presentations have uniformly close finite observations, and
corresponding atomic evaluations satisfy the required error bounds, then
every fixed continuous formula has the same controlled error.  If all errors
tend to zero with a supplied PR rate, their sentence values agree.  This is a
restricted, quantitatively witnessed form of elementary equivalence.  It
contains no assertion that every finitely satisfiable theory has a model.

\subsection{A nontrivial alternating-quantifier example}
\label{subsec:cm-minimax}

On the compact sort $[0,1]$, consider the rationally evaluated atom
\[
                  P(c,x)=|x^2-cx|
\]
and the continuous sentence
\[
             v=\inf_{c\in[0,1]}\ \sup_{x\in[0,1]} P(c,x).
\tag{M5}\label{eq:cm-example}
\]
Its quantifier order is essential: $c$ is chosen once before the adversarial
choice of $x$.  The atom takes values in $[0,1]$, and if both coordinates
change by at most $\delta$, then
\[
 |P(c,x)-P(c',x')|
 \leq |x^2-x'^2|+|cx-c'x'|\leq4\delta.
\]
Hence dyadic nets of mesh $2^{-k-2}$ give a formula-value error at most
$2^{-k}$.  Exact finite rational minimax calculations therefore define the
value in iPRA, without a real optimization oracle.

The semantic value and optimizer can also be identified explicitly.  For
$0\leq c\leq1$,
\[
       x^2-cx=(x-c/2)^2-c^2/4,
\]
so its negative part has maximum $c^2/4$, attained at $x=c/2$.
Its positive part has maximum $1-c$, attained at $x=1$: indeed
\[
       (1-c)-(x^2-cx)=(1-x)(1+x-c)\geq0.
\]
Consequently
\[
        \sup_x P(c,x)=\max\{c^2/4,1-c\}.
\]
The first expression increases and the second decreases.  Their common value
occurs at $c_*=2\sqrt2-2\in[0,1]$, and is
\[
                         v=3-2\sqrt2.
\tag{M6}\label{eq:cm-minimax-value}
\]
The following order cases may be performed first on rational $c$: its side
of the algebraic number $c_*$ is determined by the rational sign test on
$(c+2)^2-8$.  The factorization
$(c+2)^2-8=(c-c_*)(c+c_*+4)$ has a positive second factor.
The resulting non-strict estimates then extend to arbitrary named $c$ by
continuity.  Thus no exact sign test on an arbitrary real name is required.
For $c\leq c_*$ the second term is at least $1-c_*$; for $c\geq c_*$ the
first term is at least $c_*^2/4$.  This proves global minimality, and the strict
versions prove uniqueness.  The occurrence of $\sqrt2$ can be interpreted by
a fixed rational square-root name with its elementary inequalities.  All the
displayed algebraic estimates are denominator-cleared arithmetic statements.
The accompanying reference code computes the finite restricted sentence and
its minimizing net entry with exact integers and rational output.

In this example a quantitative uniqueness estimate upgrades the finite
witnesses to a coherent optimizer name.  Write
$w(c)=\max\{c^2/4,1-c\}$.  Since $c_*\geq1/2$, the two cases give
\[
                 w(c)-v\geq\tfrac14|c-c_*|.
\tag{M7}\label{eq:cm-minimax-gap}
\]
For $c\leq c_*$ use $1-c=v+c_*-c$; for $c\geq c_*$ use
$(c^2-c_*^2)/4=(c+c_*)(c-c_*)/4\geq(c-c_*)/4$.
Let $c_j$ be the first minimizing entry of the dyadic restriction with mesh
$\delta=2^{-j}$.  The generic formula bound gives $|v_j-v|\leq4\delta$,
and the inner supremum at $c_j$ differs from its finite maximum by at most
$4\delta$.  Hence $w(c_j)\leq v+8\delta$, so
$|c_j-c_*|\leq32\,2^{-j}$.  Therefore $a_n=c_{n+6}$ satisfies
$|a_n-c_*|\leq2^{-n-1}$, and
$|a_n-a_m|\leq2^{-n}$ for $m\geq n$.
The first-minimum scan thus constructs a PR fast name of the exact optimizer.
The additional ingredient is the proved gap \eqref{eq:cm-minimax-gap}, not
an assumption that arbitrary approximate minimizers are coherent.

\subsection{Sharp boundaries and the remaining Lean work}

There are three limitations with mathematical content.
First, boundedness is not enough for \eqref{eq:cm-nets}.  For example, the unit
ball of an infinite-dimensional Hilbert space has infinitely many orthogonal
unit vectors separated by a fixed positive distance.  No finite sufficiently
small net covers it.  The compact formula evaluator therefore cannot be
applied to arbitrary quantification over that unit ball.  A specified compact
subset with a uniform tail bound, a finite-rank restriction, or a different
witness contract is needed.

Second, a name of a formula value is not a decision procedure for whether it
is zero.  For a uniform example, let $h(e,t)$ be the fixed PR test that a
machine $e$ first halts at step $t\geq1$, with a harmless one-step delay
ensuring no halt at step zero.  On a one-point metric sort, give a predicate
the value $r_e=2^{-t}$ if such a first halting time exists, and $0$ otherwise.
Its $k$th rational approximation is $2^{-t}$ if a first halt is found with
$t\leq k$, and $0$ otherwise.  Bounded simulation is PR and the approximation
error is at most $2^{-k-1}$.  Its spatial modulus is trivial, and its sort has
a one-point net at every precision.  A total algorithm deciding whether
$r_e=0$ would decide nonhalting and hence halting.  Thus even this simplest
compact metric semantics cannot supply uniform exact truth tests.  This is a
metatheoretic obstruction to a proposed uniform procedure, not an
independence claim about any particular named sentence.

Third, the finite minima used above supply approximate witnesses.  They do not
supply a coherent exact witness unless an additional mechanism, such as a
contraction or a quantitative uniqueness modulus, is available.  Likewise
this construction provides no ultrafilter, saturated model, or model of an
arbitrary consistent theory.

For Lean, the principal new definitions are a fixed finite multisorted
signature, its term and continuous-formula syntax, finite environments,
rational finite evaluation, the two-sided finite covering relation, and
per-atom modulus contracts.  The first central lemma is
Theorem~\ref{thm:cm-backforth}; its proof follows the displayed syntax cases,
with the quantifier case reduced to a lemma about minima of two finite lists
related in both directions.  The second central lemma is
Theorem~\ref{thm:cm-values}, obtained by substituting two net stages and two
name stages into that result.  Approximate attainment is the first minimizing
entry from the finite scan.  No real-number library is needed for these
arithmetic endpoints; a separate semantic interpretation can identify the
result with conventional continuous logic.

The delivered \texttt{QuantitativeKernel.lean} checks a qualitative nucleus:
positive formulas built from atoms, conjunction, disjunction, and the two
quantifiers transfer across a relation with selected forward and backward
lifts, provided atoms transfer.  Its \texttt{positive\_transfer} theorem
contains no quantifier-depth loss; \texttt{eval\_correct} connects finite
Boolean evaluation to the corresponding finite-domain semantics.  These are
Lean metatheorems.  They are not a completed formalization of the numerical
continuous-logic estimates above, nor are they silently identified with
\texttt{Prf PRAAx} derivations.  To reach that endpoint, instantiate the finite
minimum lemmas in the repository's rational arithmetic, Booleanize their
bounded list invariants, and construct the syntax-indexed family of
\texttt{Prf PRAAx} proof terms.  The existing
\texttt{SetoidInterpretation.sound} provides the subsequent interpretation
interface for numerical sort families whose closure has been established.

\section{Function-space topology from explicit finite-rank operators}
\label{sec:function-topology}

The earlier topology packet already supplies covers, carriers, finite
homology calculations, and a careful warning about coherent witnesses.
The next useful step is to construct \emph{continuous functions} from
finite observations, with all denominators and continuity constants
controlled.  We give such a construction on an arbitrary presented
totally bounded metric space.  It yields a quantitative form of
Arzel\`a--Ascoli, a converse, and a gluing theorem.  These results concern
infinite spaces and function spaces even though their arithmetic proofs
use finite lists.

\subsection{The input presentation and the meaning of the statement}

Write $d_s=2^{-s}$.  A rational metric basis is a decidable collection
$A\subseteq\N$ with a fixed PR rational-valued distance $d(a,b)$ satisfying
the pseudometric laws.  Zero distance is quotiented extensionally.
For each $s$, the presentation supplies a positive natural number $N(s)$
and a list
\[
 c(s,0),\ldots,c(s,N(s)-1)\in A
\]
with the covering assertion
\begin{equation}\label{eq:top-net}
 a\in A\quad\Longrightarrow\quad
 \exists i<N(s)\;d(a,c(s,i))\le d_s.
\end{equation}
The list functions and $N$ are fixed PR programs, and the covering
assertion is an explicit premise or a previously derived theorem.
Let $X$ denote the metric completion in the external interpretation.
The list at scale $s+1$ covers $X$ at radius $d_s$ constructively:
given a regular name for $x$, take a rational approximation $a$ within
$d_{s+1}$ of $x$, and then use \eqref{eq:top-net} to find a list
member within $d_{s+1}$ of $a$.  The triangle inequality gives the
claimed bound.  This extra precision bit avoids selecting a minimum
among finitely many real distances.  All arithmetic calculations
below are performed at rational queries; extension to $X$ uses their
displayed continuity moduli.

For readers wishing to avoid any use of a real-valued supremum, every
assertion $\|f-g\|_\infty\le\eta$ below means
$\forall x\in X\;|f(x)-g(x)|\le\eta$.  It does not invoke existence of a
supremum as a new mathematical principle.

Fix a family $(f_a)_{a\in D}$ of real-valued functions on $X$, an integer
$B\ge0$, and a fixed PR modulus $\mu$ such that
\begin{equation}\label{eq:top-family}
 |f_a(x)|\le B,\qquad
 d(x,y)\le 2^{-\mu(j)}\quad\Longrightarrow\quad
 |f_a(x)-f_a(y)|\le 2^{-j}.
\end{equation}
The same modulus must work for every family member.  The symbol $a$ is
a parameter for one fixed evaluator, not a code fed to a universal PR
interpreter.  At the arithmetic level, observations of $f_a$ are
rationals $E(a,x,p)$ with error at most $d_p$ at a rational query
$x\in A$.  A fully rational specification is given below.  The family
need not have decidable equality, and its members need not be specified
by a single finite polynomial or by a discrete table.

\subsection{A denominator bound produced by the net}

\begin{lemma}[Quantitative partition of unity]\label{lem:tent-partition}
Let $r>0$ be rational and let $c_0,\ldots,c_{n-1}$ be an $r/2$-net in
$X$, with $n\ge1$.  Define
\[
 w_i(x)=\max\{0,r-d(x,c_i)\},\qquad
 W(x)=\sum_{i<n}w_i(x),\qquad
 \lambda_i(x)=\frac{w_i(x)}{W(x)}.
\]
Then $W(x)\ge r/2$, $\lambda_i(x)\ge0$, and
$\sum_{i<n}\lambda_i(x)=1$.  Moreover,
\begin{equation}\label{eq:tent-lipschitz}
 \sum_{i<n}|\lambda_i(x)-\lambda_i(y)|
 \le\frac{4n}{r}d(x,y).
\end{equation}
If $h=(h_i)_{i<n}$ and $|h_i|\le B$, the reconstructed function
\[
 R(h)(x)=\sum_{i<n}\lambda_i(x)h_i
\]
is bounded by $B$ and is Lipschitz with constant $4Bn/r$.
For two coefficient vectors,
\begin{equation}\label{eq:tent-coefficient}
 |R(h)(x)-R(\widetilde h)(x)|
 \le\max_{i<n}|h_i-\widetilde h_i|.
\end{equation}
\end{lemma}

\begin{proof}
Choose a covering index $j$ for $x$.  Then
$w_j(x)=r-d(x,c_j)\ge r/2$, and all other weights are nonnegative.
This proves the denominator bound and the first three assertions.
The triangle inequality gives
$|d(x,c_i)-d(y,c_i)|\le d(x,y)$.  The real function
$t\mapsto\max(0,r-t)$ does not increase distances, so
\[
 |w_i(x)-w_i(y)|\le d(x,y),\qquad
 |W(x)-W(y)|\le n\,d(x,y).
\]
Using positive denominators, add and subtract $w_i(y)/W(x)$ to obtain
\[
 |\lambda_i(x)-\lambda_i(y)|
 \le\frac{|w_i(x)-w_i(y)|}{W(x)}
  +\frac{w_i(y)|W(y)-W(x)|}{W(x)W(y)}.
\]
Sum this inequality.  In the second term the sum of $w_i(y)$ is
$W(y)$, and hence
\[
 \sum_i|\lambda_i(x)-\lambda_i(y)|
 \le\frac{n\,d(x,y)+|W(y)-W(x)|}{W(x)}
 \le\frac{4n}{r}d(x,y).
\]
Multiply by the coefficient bound $B$ to obtain the Lipschitz assertion.
The bound $|R(h)(x)|\le B$ and \eqref{eq:tent-coefficient} follow by
nonnegativity of the weights and their sum being one.  No cancellation,
choice of a nearest centre, or decision about equality of real points
is used.
\end{proof}

The formula is a continuous partition of unity; smoothness is not
claimed.  On the rational basis, all weights and all denominators are
exact rational numbers.  The inverse operation is justified by a
\emph{derived} lower bound $r/2$, rather than by a new oracle promising
that a denominator is nonzero.

\subsection{An explicit quantitative Arzel\`a--Ascoli theorem}

\begin{theorem}[Finite continuous approximants with an entropy bound]
\label{thm:quant-ascoli}
Under \eqref{eq:top-net} and \eqref{eq:top-family}, put
\begin{align*}
 s_k&=\mu(k+2)+2,& r_k&=2^{-\mu(k+2)},& n_k&=N(s_k),\\
 \delta_k&=2^{-(k+3)},&
 M_k&=2B\,2^{k+3}+1,& H_k&=M_k^{n_k}.
\end{align*}
The basis list has mesh $r_k/4$ and hence covers $X$ with mesh $r_k/2$.
Using the centres $c(s_k,i)$, form the weights of
Lemma~\ref{lem:tent-partition} with radius $r_k$.
There is an explicit list of $H_k$ continuous functions $g_{k,h}$
such that
\begin{enumerate}
 \item $|g_{k,h}(x)|\le B$ and each $g_{k,h}$ is Lipschitz with
 constant at most $4Bn_k/r_k$;
 \item for every $a\in D$, finitely many evaluations of $f_a$ at the
 $n_k$ net points, each to accuracy $\delta_k$, produce an index
 $h<H_k$ satisfying
 \begin{equation}\label{eq:ascoli-cover}
  \|f_a-g_{k,h}\|_\infty\le 2^{-(k+1)};
 \end{equation}
 \item the list size, its rational evaluators, the selected index, and
 every bound above are PR in $k,B$ and the numerical presentation data,
 for each fixed collection of input programs.
\end{enumerate}
The centres of this cover lie in the ambient space of continuous
functions.  They are not asserted to belong to the original family.
\end{theorem}

\begin{proof}
The grid
\[
 G_k=\{-B,-B+\delta_k,\ldots,B-\delta_k,B\}
\]
has exactly $M_k$ members because $B$ is an integer and
$\delta_k^{-1}=2^{k+3}$.  If $B=0$, the single vector is the zero
vector.  Otherwise enumerate all vectors in $G_k^{n_k}$ by
base-$M_k$ digits.  For such a vector $h$, define
$g_{k,h}=R(h)$.  Lemma~\ref{lem:tent-partition} gives the first
assertion.  Each list element has a rational evaluator on $A$ given by
the displayed finite sum and quotient; the Lipschitz bound extends it
uniquely to the completion.

For a particular $f_a$, obtain a rational $e_i$ satisfying
$|e_i-f_a(c(s_k,i))|\le\delta_k$.  Define
\[
 h_i=\operatorname{clip}_{[-B,B]}
       \bigl(\delta_k\lfloor e_i/\delta_k\rfloor\bigr).
\]
This is a decidable rational operation, including at grid boundaries.
Projection onto an interval does not increase distance from a point of
that interval.  Rounding down changes a rational by at most $\delta_k$.
Since $f_a(c(s_k,i))\in[-B,B]$, it follows that
\begin{equation}\label{eq:ascoli-sample-error}
 |h_i-f_a(c(s_k,i))|\le 2\delta_k=2^{-(k+2)}.
\end{equation}
The resulting vector lies in $G_k^{n_k}$, so its base-$M_k$ code is
less than $H_k$.

First fix $x\in A$, so that every support comparison is decidable.
If $w_i(x)>0$, then
$d(x,c(s_k,i))<r_k=2^{-\mu(k+2)}$.  The common modulus therefore gives
\[
 |f_a(c(s_k,i))-f_a(x)|\le 2^{-(k+2)}.
\]
For these indices, \eqref{eq:ascoli-sample-error} implies
$|h_i-f_a(x)|\le 2^{-(k+1)}$.  Indices with zero weight contribute
zero.  Taking the weighted average yields
\[
 |g_{k,h}(x)-f_a(x)|
 \le\sum_i\lambda_i(x)|h_i-f_a(x)|
 \le 2^{-(k+1)}\sum_i\lambda_i(x)=2^{-(k+1)}.
\]
This estimate holds for every rational query.  For a regular name of
an arbitrary $x\in X$, the supplied modulus for $f_a$ and the proved
Lipschitz bound for $g_{k,h}$ pass the same non-strict estimate to $x$.
No sign decision about a real distance is needed.

For the PR claim, exponentiation, finite list formation, base expansion,
signed rational floors, clipping, and bounded sums are PR.  The call
depth of the input evaluator is fixed; only its numerical arguments
vary.  Every loop is bounded by a displayed natural number, although
$H_k$ can be very large.  In particular no test of the form ``does this
function agree with some earlier function?'' is part of the enumeration.
\end{proof}

\begin{remark}[What the theorem supplies]
This is a uniform finite cover in the supremum-distance sense.  If one
uses a definition of total boundedness demanding that the centres
belong to the given family, additional witnesses for the inhabited
profile classes are needed.  Given one member of the family in each
inhabited profile class, the triangle inequality changes the radius
from $2^{-(k+1)}$ to $2^{-k}$.  The displayed construction itself does
not decide which profile classes are inhabited.
Likewise, replacing the supplied common modulus by bare pointwise
equicontinuity is an extra theorem; it is not a consequence of the
finite enumeration.
\end{remark}

\begin{corollary}[Finite-rank approximation of continuous functions]
\label{cor:finite-rank}
With the same weights, let
\[
 (P_k f)(x)=\sum_{i<n_k}\lambda_i(x)f(c(s_k,i)).
\]
Then $P_k$ is a positive linear operator, preserves constants, has
operator norm at most one, and its image is contained in the span of
$\lambda_0,\ldots,\lambda_{n_k-1}$.  Every function satisfying the
given modulus obeys
\[
 \|P_k f-f\|_\infty\le 2^{-(k+2)}.
\]
Thus the metric presentation produces finite-rank approximation
operators with an explicit convergence rate on every family with a
common modulus.
\end{corollary}

\begin{proof}
Linearity is distributivity of a finite sum.  Positivity, preservation
of constants, and the norm bound follow from the nonnegative weights
adding to one.  The range statement is immediate from the formula.
At a rational query $x$, every nonzero weight comes from a centre within
$2^{-\mu(k+2)}$ of $x$.  Apply the modulus to each such summand and
average, exactly as in the preceding proof but without quantization;
then extend the estimate along regular names.
\end{proof}

The operator statement uses ordinary continuous-function language in
the external interpretation.  Its arithmetic content is the family
of finite sum identities and error estimates.  It does not assert that
all continuous functions form a numerical sort of \ipra.

For example, on $[0,1]^d$ with the maximum metric, the dyadic grid
gives $N(s)=(2^s+1)^d$.  Functions bounded by the integer $B$ and
Lipschitz with constant at most $2^\ell$ have
$\mu(j)=j+\ell$.  The theorem therefore gives the entirely explicit
covering bound
\[
 H_k=(2B\,2^{k+3}+1)^{(2^{k+\ell+4}+1)^d}.
\]
The estimate is deliberately conservative.  Its purpose is to expose
the full primitive-recursive dependence and provide a reusable proof,
not to claim optimal metric entropy.

\subsection{The converse: finite approximants recover regularity}

\begin{proposition}[A quantitative converse]\label{prop:ascoli-converse}
Suppose that for every $j$ a finite nonempty list
$g_{j,0},\ldots,g_{j,L_j-1}$ covers a family $\mathcal F$ within
supremum distance $2^{-j}$.  Suppose each $g_{j,i}$ has an explicit
uniform modulus $\nu(j,i,t)$, and the list at $j=0$ has integer bounds
$|g_{0,i}|\le C_i$.  Then all $f\in\mathcal F$ satisfy
\[
 |f|\le1+\max_{i<L_0}C_i
\]
and share the modulus
\begin{equation}\label{eq:ascoli-converse-modulus}
 \mu_*(k)=\max_{i<L_{k+2}}\nu(k+2,i,k+1).
\end{equation}
If the displayed data are given by fixed PR programs, these new data
are PR as well.  It is enough to have the finite-cover assertion;
an algorithm choosing a covering index is not required for this
direction.
\end{proposition}

\begin{proof}
For the range bound, take a covering index for $f$ in the list at
$j=0$ and use the triangle inequality.  For the modulus, take a
covering index $i$ at $j=k+2$.  If
$d(x,y)\le2^{-\mu_*(k)}$, it is at most
$2^{-\nu(k+2,i,k+1)}$.  Consequently
\begin{align*}
 |f(x)-f(y)|
 &\le |f(x)-g_{k+2,i}(x)|
   +|g_{k+2,i}(x)-g_{k+2,i}(y)|
   +|g_{k+2,i}(y)-f(y)|\\
 &\le 2^{-(k+2)}+2^{-(k+1)}+2^{-(k+2)}=2^{-k}.
\end{align*}
The index is eliminated from a finite existential hypothesis after
this estimate; no countable choice is used.  Maxima of finite lists
are computed by a bounded recursion.
\end{proof}

The forward and reverse results identify the effective content of
Ascoli approximation.  Uniform range and continuity information
produce finite approximants, and finite approximants with their own
regularity information recover a common modulus.  Convergence of an
arbitrary sequence is a separate question.

\subsection{Gluing guarded local data}

\begin{theorem}[Continuous gluing with a produced modulus]
\label{thm:guarded-gluing}
Let $c_0,\ldots,c_{n-1}$ be an $r/2$-net.  For each $i$, suppose a
local function $f_i$ is given on the open ball
$\{x:d(x,c_i)<2r\}$.  Assume that on this ball $|f_i|\le B$ and
$f_i$ is Lipschitz with constant $L$.  Assume also that on every
overlap where both distances are less than $r$,
\[
 |f_i(x)-f_j(x)|\le\alpha.
\]
On the rational basis define $p_i(x)=w_i(x)f_i(x)$ when
$d(x,c_i)<r$ and $p_i(x)=0$ otherwise; these are decidable rational
support tests.  Here $w_i$ is as in Lemma~\ref{lem:tent-partition}.
The products have unique Lipschitz extensions to $X$.  For those
extensions,
\[
 g(x)=\frac{\sum_{i<n}p_i(x)}{W(x)}
\]
is a globally defined function with $|g|\le B$,
\begin{equation}\label{eq:gluing-lip}
 |g(x)-g(y)|\le
 \left(2nL+\frac{4nB}{r}\right)d(x,y),
\end{equation}
and $|g(x)-f_j(x)|\le\alpha$ whenever $d(x,c_j)<r$.
In particular exact overlap agreement yields a genuine global
continuous gluing of the restrictions to the $r$-balls.  The assertions
also hold for complex values, with
the absolute value replaced by the complex norm.
\end{theorem}

\begin{proof}
First take $x,y$ in the rational basis.  The function $p_i$ is
well-defined because the support where its
formula uses $f_i$ lies strictly inside the domain of $f_i$.
If $x$ and $y$ both lie in this support, then
\begin{align*}
 |p_i(x)-p_i(y)|
 &\le w_i(x)|f_i(x)-f_i(y)|
   +|w_i(x)-w_i(y)|\,|f_i(y)|\\
 &\le (rL+B)d(x,y).
\end{align*}
If neither lies in the support, both products vanish.  If only $x$
does, the reverse triangle inequality gives
$w_i(x)=r-d(x,c_i)\le d(x,y)$ because $d(y,c_i)\ge r$.
Thus $|p_i(x)|\le B\,d(x,y)$, which is also bounded by
$(rL+B)d(x,y)$.  The case with only $y$ is symmetric.  Therefore every
$p_i$ extends continuously across the support boundary and has the
stated global Lipschitz bound.

Put $P=\sum_i p_i$.  We have
\[
 |P(x)-P(y)|\le n(rL+B)d(x,y),\qquad |P(y)|\le B W(y).
\]
Use $W(x)\ge r/2$ and $|W(x)-W(y)|\le n\,d(x,y)$ to estimate
\begin{align*}
 \left|\frac{P(x)}{W(x)}-\frac{P(y)}{W(y)}\right|
 &\le \frac{|P(x)-P(y)|}{W(x)}
   +\frac{|P(y)|\,|W(y)-W(x)|}{W(x)W(y)}\\
 &\le\frac{2n(rL+2B)}{r}d(x,y),
\end{align*}
which is \eqref{eq:gluing-lip}.  These rational-query estimates give
unique extensions to $X$ and remain true there by passage along
regular names with the displayed modulus.  The range bound follows
from $|P(x)|\le BW(x)$.  Finally, at a rational query with
$d(x,c_j)<r$, every index contributing
to the weighted average also has $d(x,c_i)<r$.  The overlap hypothesis
therefore gives
\[
 |g(x)-f_j(x)|
 \le\sum_i\lambda_i(x)|f_i(x)-f_j(x)|\le\alpha.
\]
For an arbitrary $x\in X$ with $d(x,c_j)<r$, sufficiently accurate
rational approximations lie in the same guarded ball; the Lipschitz
bounds for $g$ and $f_j$ transfer this estimate to $x$.
Only positivity, the triangle inequality, and finite sums were used,
so the proof is unchanged for complex values.
\end{proof}

For an implementation on the rational basis, the support test compares
two rationals.  The uniform arithmetic theorem uses one fixed total
PR program $E(i,x,t)$ for the patch parameter $i<n$, whose correctness
is required only inside the corresponding guard.  An externally
fixed finite collection of local programs may instead be compiled into
such an evaluator.  This does not permit unbounded interpretation of
arbitrary input program codes.  The program never
uses a local value outside the support of its weight.  The larger
radius $2r$ supplies room for evaluating names near support boundaries;
the preceding global modulus then extends the rational construction to
$X$.  For real queries one uses this extension, not an exact decision
whether $d(x,c_i)=r$.

More explicitly, if $E_i(x,t)$ approximates $f_i(x)$ within $2^{-t}$
on its guard, evaluate
\[
 G(x,t)=\frac{\sum_{i:\,d(x,c_i)<r}w_i(x)E_i(x,t)}{W(x)}.
\]
Every rational query to a local evaluator is within its guard, and
$|G(x,t)-g(x)|\le2^{-t}$ because the nonnegative normalized weights
sum to one.  Thus the gluing construction also supplies an explicit
regular-name evaluator, not just a semantic existence assertion.

This solves a specific presentation problem: local continuous data
with a certified finite metric cover can be glued by a concrete
formula, and the reciprocal bound and continuity modulus are computed
from the cover.  The construction does not manufacture finite covers
from an arbitrary open-cover enumeration, or smooth local data from
mere continuous data.

\subsection{Why regular subsequence extraction remains unavailable}

The distinction between approximation and subsequence selection is
already visible for constant functions.  Here is a self-contained
computability obstruction of the type introduced by Specker
\cite{Specker1949}; it also explains why one should keep the different
Ascoli principles separate \cite{Kreuzer2014}.

Fix a bounded PR simulator for Turing machines and set
\[
 q_s=\sum_{\substack{e<s\\ \text{machine }e\text{ on }e
                      \text{ halts by step }s}}4^{-(e+1)}.
\]
This is a PR sequence of rational numbers, is increasing, and lies in
$[0,1/3]$.  In the external classical interpretation its limit is
\[
 q=\sum_{e\in K}4^{-(e+1)},
\]
where $K$ is the diagonal halting set.  This real is not computable.
Indeed, after the bits of $K$ below $e$ have been determined, subtract
their contributions.  If $e\notin K$, the remaining value is at most
$4^{-(e+1)}/3$, whereas if $e\in K$ it is at least $4^{-(e+1)}$.
An approximation with error less than $4^{-(e+1)}/6$ distinguishes
these cases by comparison with $2\,4^{-(e+1)}/3$.  A computable name
for $q$ would consequently decide $K$, a contradiction.

Now let $f_s(x)=q_s$ on $[0,1]$.  The family has bound $B=1$ and the
common modulus $\mu(j)=0$ for every $j$.  Its finite approximants are
particularly easy.  Nevertheless, a computable strictly increasing
subsequence together with a computable modulus of uniform convergence
would compute $q$: evaluate a constant subsequence member at the index
specified by the modulus.  Thus no such regular subsequence can be
produced in general.  This is an external impossibility theorem about
uniform computations, not a new internally formalized incompleteness
claim for \ipra.  It does not rule out a particular sequence having a
supplied modulus or a problem-specific compactness argument producing
one.

\subsection{The exact rational target and a Lean development path}

For implementation, the semantic functions $f_a$ can be replaced
throughout by the following rational contracts, where $x,y\in A$ and
$p,q,j\in\N$:
\begin{align}
 |E(a,x,p)-E(a,x,q)|&\le d_p+d_q,\label{eq:top-name}\\
 |E(a,x,p)|&\le B+d_p,\label{eq:top-range}\\
 d(x,y)\le2^{-\mu(j)}\quad&\Longrightarrow\quad
 |E(a,x,p)-E(a,y,q)|\le d_j+d_p+d_q.
 \label{eq:top-space}
\end{align}
Together with the metric and net laws, these are sufficient for the
finite approximation theorem.  They state precisely what a rational
evaluation interface must provide.  They are hypotheses about fixed
programs, not automatically true properties of arbitrary program
codes.

At precision $k$, use $e_i=E(a,c(s_k,i),k+3)$ and the clipped floor
defined above.  Since \eqref{eq:top-range} gives
$e_i\in[-B-\delta_k,B+\delta_k]$, a direct three-case comparison
shows
\[
 |h_i-e_i|\le\delta_k.
\]
For $e_i\in[-B,B]$ this is the rounding bound.  For $e_i>B$ or
$e_i<-B$, clipping returns the respective endpoint and the range
contract gives the same bound.  If $w_i(x)>0$,
\eqref{eq:top-space} with $j=k+2$ now gives
\[
 |h_i-E(a,x,p)|
 \le\delta_k+d_{k+2}+d_{k+3}+d_p
 =d_{k+1}+d_p.
\]
Averaging proves the exact first-order target
\begin{equation}\label{eq:top-object-target}
 \boxed{\quad
 |g_k(a,x)-E(a,x,p)|\le 2^{-(k+1)}+2^{-p}.
 \quad}
\end{equation}
Every term in this formula is a rational code.  After checking
positive denominators, multiplication clears them, yielding a
quantifier-free integer inequality.  The net witness needed for the
denominator is an existential premise used outside the induction over
the finite sum.  In particular there is no induction on a statement
quantifying over all functions, all real points, or all precisions.

A Lean implementation should use two deliberately separate layers.
The following list gives the definitions and theorem order; the names
are proposed module names, not claims about compiled code delivered
elsewhere in this repository.
\begin{enumerate}
 \item \texttt{RationalMetricBasis} contains a decidable domain, a
 rational distance program, the pseudometric laws, \texttt{netSize},
 \texttt{netPoint}, and \eqref{eq:top-net}.  The arithmetic version
 stores PR syntax and object derivations for those laws; an ordinary
 Lean record with rational functions is the semantic specification.
 \item \texttt{TentWeights} defines the rational positive part,
 weights, denominator, and coefficient reconstruction.  Prove
 \texttt{weight\_nonnegative}, \texttt{denominator\_lower},
 \texttt{weight\_difference}, \texttt{partition\_variation}, and
 \texttt{reconstruction\_nonexpansive}, in that order.  Finite-sum
 estimates use existing sum linearity and rational-order derivations.
 \item \texttt{EquicontinuousEvaluator} stores $B,\mu,E$ and the
 three contracts \eqref{eq:top-name}--\eqref{eq:top-space}.  Add
 \texttt{clipFloor}, its grid-membership and error lemmas, and
 mixed-radix encoding of coefficient lists.
 \item \texttt{AscoliApproximation} defines $s_k,r_k,n_k,M_k,H_k$
 and $g_k$ literally as above.  Derive
 \texttt{ascoli\_finite\_error} with target
 \eqref{eq:top-object-target}, and separately the grid-code bound.
 No existential over function names is needed: $g_k$ is an explicit
 rational expression from the given observations.
 \item \texttt{GuardedGluing} handles the two support cases for each
 product, proves its Lipschitz bound, then sums and divides using
 \texttt{denominator\_lower}.  A semantic completion theorem
 identifies the resulting regular name with the continuous gluing.
 \item \texttt{AscoliConverse} introduces a bounded covering index,
 uses the three-term triangle estimate, and removes the index by
 existential elimination.  The maximum of a finite list supplies
 \eqref{eq:ascoli-converse-modulus}.
\end{enumerate}

For the object endpoint, the goal is an actual derivation
$\Gamma\vdash[\mathrm{PRAAx}]\varphi$ after inserting the chosen
PR metric, net, and evaluator terms.  The context $\Gamma$ retains
their explicit contracts.  A semantic Lean proof about
\texttt{MetricSpace} or \texttt{ContinuousMap} is useful for the
external interpretation but cannot by itself discharge this endpoint.
The required arithmetic induction invariants concern finite partial
sums and bounded coefficient lists and therefore fit the existing
quantifier-free derivation discipline.

\section{Coercive equations in presented Hilbert spaces}
\label{sec:hilbert}

This section supplies an infinite-dimensional application of the framework.
Its constructive content is stronger than verifying a proposed solution:
the input data produce a regular solution name with a prescribed
primitive-recursive clock. It also identifies why the corresponding statement
for an arbitrary bounded bilinear form and a weakly presented functional
cannot be asserted without additional information. Computable versions of
Riesz representation and Lax--Milgram have been studied previously
\cite{BrattkaYoshikawa2006}; the argument here is a self-contained quantitative
specialization designed for the arithmetic target.

\subsection{The presentation and its actual inputs}

\begin{definition}[Rational pre-Hilbert presentation]
A rational pre-Hilbert presentation is a coded rational vector space \(H_0\)
with primitive-recursive zero, addition and rational scalar multiplication,
and a rational-valued symmetric positive-definite inner product
\(\langle\ ,\ \rangle\). All vector-space and inner-product laws have their
ordinary finite arithmetic proofs. Its named completion \(H\) consists of
regular names, identified by distance zero. We write
\(\|v\|^2=\langle v,v\rangle\).
\end{definition}

Examples are finite rational sequences with their Euclidean inner product,
or rational polynomials vanishing at \(0,1\), with inner product
\(\int_0^1p'q'\). The dimension of a finite code may grow during an
algorithm. There is no hypothesis that \(H\) is finite-dimensional or totally
bounded.

Norms need not be computed by a real square root. For \(v\in H_0\) and
rational \(r\geq0\), the test \(\|v\|\leq r\) means
\(\langle v,v\rangle\leq r^2\). Rational upper approximations to a norm can
be obtained by a bounded integer square-root calculation. For example,
search integers \(j\) up to
\(\lceil2^t(2+\langle v,v\rangle)\rceil\) for the first
\((j2^{-t})^2\geq\langle v,v\rangle\).
The returned upper radius differs from the norm by at most \(2^{-t}\).

\begin{definition}[Strong operator and vector data]
\label{def:strong-operator}
Fix primitive-recursive maps
\[
 a:H_0\times\N\longrightarrow H_0,\qquad b:\N\longrightarrow H_0.
\]
They name a real-linear operator \(A:H\to H\) and a vector \(b_\infty\in H\)
with errors
\[
 \|a(v,k)-Av\|\leq2^{-k},\qquad
 \|b(k)-b_\infty\|\leq2^{-k}.
 \tag{H1}
\]
Positive rational inputs \(m,M\), with \(m\leq M\), and a rational \(B\geq0\)
certify
\[
 \langle Av,v\rangle\geq m\|v\|^2,\qquad
 \|Av\|\leq M\|v\|,\qquad \|b_\infty\|\leq B.
 \tag{H2}
\]
The maps and their proved certificates, not an unnamed abstract operator, are
the inputs to the arithmetic construction.
\end{definition}

The notation in (H1)--(H2) has a finite reading. One supplies the pairwise
Cauchy estimates for \(a(v,k)\) and \(b(k)\); approximate linearity, for example
\[
\|a(sv+tw,k)-s\,a(v,k)-t\,a(w,k)\|
       \leq(1+|s|+|t|)2^{-k};
\]
and boundedness and coercivity estimates with approximation errors.
For coercivity an adequate entirely rational contract is
\[
\langle a(v,k),v\rangle
 \geq m\langle v,v\rangle-2^{-k}(1+\langle v,v\rangle).
\]
For boundedness one may use every rational upper radius \(r\) of \(v\):
\[
\langle v,v\rangle\leq r^2
 \ \Longrightarrow\ \|a(v,k)\|\leq Mr+2^{-k}.
\]
Here \(r\) is a free numerical parameter. These contracts define the named
operator by extension from the dense space and prove (H2); rational
separation removes auxiliary errors. Conversely, (H1)--(H2) immediately
give these finite contracts. For the exact \(2^{-k}\) approximation in (H1),
one can normalize the supplied Cauchy family by a fixed precision shift.
One must prove the contracts for a concrete presentation; merely storing an
assertion alongside the code is insufficient.

\subsection{The elementary quantitative engine}

\begin{lemma}[Perturbed contraction]
\label{lem:perturb}
Let \(0\leq q<1\), \(\delta=1-q\), and suppose nonnegative quantities satisfy
\[
             e_{j+1}\leq q e_j+\eta,\qquad e_0\leq R .
\]
Then
\[
 e_j\leq q^jR+\eta\sum_{i<j}q^i
       =q^jR+\frac{\eta(1-q^j)}{\delta}
       \leq q^jR+\frac{\eta}{\delta}.
 \tag{H3}
\]
All claims about rational inputs are quantifier-free induction consequences.
\end{lemma}
\begin{proof}
At \(j=0\) the finite sum is zero. Multiplying the inductive estimate by
\(q\geq0\), adding \(\eta\), and reindexing gives the next
estimate. The geometric identity follows by subtracting \(q\) times the
sum from the sum. The denominator is positive. These are finite addition,
multiplication, comparison and induction steps.
\end{proof}

\begin{lemma}[A clock without logarithms]
\label{lem:clock}
If \(0<\delta\leq1\), \(q=1-\delta\), and \(h\in\N\) satisfies \(h\delta\geq1\),
then \(q^h\leq\tfrac12\). Consequently, for integers \(R\geq1\),
\[
       N_n=h(n+R+3)\quad\Longrightarrow\quad
       Rq^{N_n}\leq2^{-n-3}.
 \tag{H4}
\]
\end{lemma}
\begin{proof}
Induction gives
\((1+j\delta)(1-\delta)^j\leq1\).
For the step use
\[
(1+(j+1)\delta)(1-\delta)
=1+j\delta-(j+1)\delta^2\leq1+j\delta.
\]
At \(j=h\), \(1+h\delta\geq2\).
Now \(q^{h(n+R+3)}\leq2^{-(n+R+3)}\).
The elementary inequality \(R\leq2^R\) proves (H4).
\end{proof}

\begin{lemma}[Coercivity supplies the contraction gap]
\label{lem:coercive-contraction}
For the data in Definition~\ref{def:strong-operator}, put
\[
 \lambda=\frac{m}{M^2},\quad
 \delta=\frac{m^2}{2M^2},\quad q=1-\delta,\quad
 T(v)=v-\lambda Av+\lambda b_\infty .
 \tag{H5}
\]
Then \(1/2\leq q<1\) and
\[
                     \|T(v)-T(w)\|\leq q\|v-w\|.
 \tag{H6}
\]
\end{lemma}
\begin{proof}
Write \(z=v-w\). Real linearity and (H2) give
\[
\begin{split}
\|z-\lambda Az\|^2
 &=\|z\|^2-2\lambda\langle Az,z\rangle+\lambda^2\|Az\|^2\\
 &\leq(1-2\lambda m+\lambda^2M^2)\|z\|^2\\
 &=(1-2\delta)\|z\|^2
 \leq(1-\delta)^2\|z\|^2 .
\end{split}
\]
Both comparison radii are nonnegative, so this proves (H6).
The finite rational proof expands the same squared inner product. Auxiliary
precision parameters account for evaluations of \(A\), and their errors are
removed by Lemma~\ref{lem:separation}. No spectral theorem, eigenvector
selection, orthonormal basis, or minimization principle is required.
\end{proof}

\subsection{A solution producer with all resource bounds}

Choose the explicit natural numbers
\[
\begin{split}
h&=\ceil{1/\delta},\qquad
R=\max(1,\ceil{\lambda B/\delta}),\qquad
C=\max(1,\ceil{2\lambda/\delta}),\\
k_n&=n+3+C,\qquad N_n=h(n+R+3),\qquad
\eta_n=\delta\,2^{-n-3}.
\end{split}
\tag{H7}
\]
For requested output precision \(n\), run the finite recurrence
\[
\begin{split}
 z^{\,n}_0&=0,\\
 z^{\,n}_{j+1}
   &=z^{\,n}_j-\lambda a(z^{\,n}_j,k_n)+\lambda b(k_n)
       \quad(j<N_n),\\
 c_n&=z^{\,n}_{N_n}.
\end{split}
\tag{H8}
\]
Every operation in (H7)--(H8) is primitive recursive on finite codes.
On malformed inputs or invalid positive constants the producer can
return \(0\); correctness is asserted on the valid domain.

\begin{theorem}[Quantitative operator Lax--Milgram]
\label{thm:hilbert-solver}
The sequence \(c_n\) produced by (H8) is regular and defines a unique
solution \(u\) of \(Au=b_\infty\). More precisely,
\[
\begin{gathered}
       \|c_n-c_\ell\|\leq2^{-n-1}\quad(\ell\geq n),\\
       \|c_n-u\|\leq2^{-n-2},\qquad
       \|u\|\leq B/m.
\end{gathered}
\tag{H9}
\]
The construction is uniform in the numerical inputs and schematic in the
fixed primitive-recursive operator and vector producers.
\end{theorem}
\begin{proof}
Let \(F_n(v)\) denote the finite update in (H8). Since \(2^C\geq C\),
\[
 \|F_n(v)-T(v)\|\leq2\lambda\,2^{-k_n}
                  \leq\delta2^{-n-3}=\eta_n.
 \tag{H10}
\]
Thus for any precision levels and inputs,
\[
\|F_n(v)-F_\ell(w)\|
 \leq q\|v-w\|+\eta_n+\eta_\ell .
\tag{H11}
\]
This is the central finite comparison contract.

First bound every intermediate state. Since
\(\|T(0)\|\leq\lambda B\leq\delta R\), induction gives
\[
\|z^{\,\ell}_j\|\leq R+\eta_\ell/\delta.
\tag{H12}
\]
Multiplication by \(q\), followed by addition of
\(\delta R+\eta_\ell\), preserves this bound.

Now fix \(\ell\geq n\), so \(N_\ell\geq N_n\).
Compare the \(n\)-run with the \emph{last \(N_n\) steps} of the \(\ell\)-run:
\[
             z^{\,n}_i
       \quad\hbox{and}\quad z^{\,\ell}_{N_\ell-N_n+i},
                    \qquad 0\leq i\leq N_n.
\]
Their initial distance is at most \(R+\eta_\ell/\delta\) by (H12).
Lemma~\ref{lem:perturb} and (H11) bound the final distance by
\[
\begin{split}
q^{N_n}(R+\eta_\ell/\delta)+(\eta_n+\eta_\ell)/\delta
 &\leq Rq^{N_n}+\eta_n/\delta+2\eta_\ell/\delta\\
 &\leq2^{-n-3}+2^{-n-3}+2\cdot2^{-\ell-3}\\
 &\leq2^{-n-1}.
\end{split}
\tag{H13}
\]
This proves regularity directly from two finite executions. The existence of
\(u\) was not assumed to prove convergence. The completion interpretation
of this regular name gives \(u\).
Taking arbitrarily fine comparison indices \(\ell\) in (H13) gives the sharper
\(\|c_n-u\|\leq Rq^{N_n}+\eta_n/\delta\leq2^{-n-2}\).
In the arithmetic reading this is a family of rational inequalities and
rational separation, not an added real limit rule.

To prove the equation using only comparisons of finite states, extend a run
by one more step and set \(w_j=\|z^{\,n}_{j+1}-z^{\,n}_j\|\).
By (H11),
\[
w_0\leq\lambda B+\eta_n,\qquad
w_{j+1}\leq q w_j+2\eta_n.
\]
Consequently
\[
\begin{split}
\|T(c_n)-c_n\|
 &\leq w_{N_n}+\eta_n\\
 &\leq q^{N_n}(\lambda B+\eta_n)+2\eta_n/\delta+\eta_n\\
 &\leq(2+3\delta)\,2^{-n-3}.
\end{split}
\tag{H14}
\]
The last bound is deliberately loose and tends to zero with an explicit
dyadic rate. Continuity of \(T\), the approximation estimate, and (H14)
show \(T(u)=u\). Since \(\lambda>0\), this is \(Au=b_\infty\).

If \(Au=Av\), coercivity applied to \(u-v\) gives
\(m\|u-v\|^2\leq0\), hence equality of the represented points.
Finally
\[
m\|u\|^2\leq\langle Au,u\rangle
 =\langle b_\infty,u\rangle\leq B\|u\|
\]
implies \(\|u\|\leq B/m\). This implication can be checked by rational upper
and lower tests; it does not require deciding whether the real norm is zero.
\end{proof}

\paragraph{Why this is an arithmetic proof plan.}
The state bound (H12), shifted two-run bound (H13), and adjacent-state bound
in (H14) are finite loop invariants. After replacing norms by rational upper
radii they are quantifier-free inequalities in state codes, loop counters,
constants, and chosen precision indices. A final observation point or test
vector is fixed before induction. The exact \(T\) in the notation is
eliminated using finite evaluators at an auxiliary precision; all additional
errors have explicit finite geometric-sum bounds. Separation yields (H11).
There is no induction on the formula ``\(z_j\) denotes a solution'' or on a
formula quantifying over the Hilbert space.
Section~\ref{sec:formalization} specifies the object-code obligations.

\begin{corollary}[Residual certification and stability]
\label{cor:residual}
For every represented \(v\),
\[
                   \|v-u\|\leq \frac{\|Av-b_\infty\|}{m}.
\tag{H15}
\]
Suppose a second operator \(\widetilde A\) has the same coercivity lower
bound \(m\), solution \(\widetilde u\), and data bound
\(\|\widetilde b\|\leq\widetilde B\). If
\[
\|b_\infty-\widetilde b\|\leq\beta,\qquad
\|(A-\widetilde A)x\|\leq\xi\|x\|,
\]
then
\[
       \|u-\widetilde u\|
       \leq \frac{\beta}{m}+\frac{\xi\widetilde B}{m^2}.
\tag{H16}
\]
\end{corollary}
\begin{proof}
For \(w=v-u\), coercivity and Cauchy--Schwarz give
\(m\|w\|^2\leq\langle Av-b_\infty,w\rangle
\leq\|Av-b_\infty\|\|w\|\).
The same nonnegative inequality argument as above proves (H15).
Apply it to \(\widetilde u\), whose residual under \(A\) is
\((A-\widetilde A)\widetilde u+\widetilde b-b_\infty\), and use
\(\|\widetilde u\|\leq\widetilde B/m\).
\end{proof}

Solving is therefore uniformly continuous on a stratum with a supplied
positive coercivity bound. The constants explain precisely why no uniform
extension to \(m=0\) is asserted.

\subsection{A real boundary-value problem, including variable coefficients}
\label{subsec:pde}

The following application constructs the strong operator instead of
postulating it. Let \(P=\Q[x]\) and
\(P_0=\{p\in\Q[x]:p(0)=p(1)=0\}\).
Equip \(P\) with inner product \(\int_0^1pq\), and \(P_0\) with
\[
                 \langle p,q\rangle_E=\int_0^1p'(x)q'(x)\,dx .
\]
All integrals are rational, since
\(\int_0^1x^j\,dx=1/(j+1)\).
These forms are positive definite. For a nonzero rational polynomial of degree
\(d\), one of \(d+1\) distinct rational interior points has nonzero value:
otherwise its coefficient vector is annihilated by an invertible Vandermonde
matrix. A coefficient bound gives an explicit Lipschitz constant, and hence
a rational interval on which its square has a positive rational lower bound.
Its squared integral is therefore positive. For the energy form, a polynomial
with zero derivative is constant, and the endpoint condition forces that
constant to be zero. This also gives finite arithmetic proofs of the
pre-Hilbert premises.
Let \(L^2_P\) and \(H_E\) denote the named completions.
These are the chosen presentations of the \(L^2\) and zero-boundary energy
spaces; no general measurable-set universe is needed.

\begin{lemma}[Polynomial inequalities and Green operator]
\label{lem:green}
For \(v\in P_0\),
\[
          \sup_{0\leq x\leq1}|v(x)|\leq\|v\|_E,\qquad
          \|v\|_{L^2}\leq\|v\|_E .
\tag{H17}
\]
For \(p\in P\) define
\[
(Gp)(x)=x\int_0^1(1-t)p(t)\,dt-\int_0^x(x-t)p(t)\,dt .
\tag{H18}
\]
Then \(Gp\in P_0\), \(-(Gp)''=p\), and
\[
  \langle Gp,v\rangle_E=\int_0^1p v,\qquad
  \|Gp\|_E\leq\|p\|_{L^2}.
\tag{H19}
\]
\end{lemma}
\begin{proof}
The finite polynomial identity \(v(x)=\int_0^x v'(t)\,dt\) and
Cauchy--Schwarz give
\(|v(x)|^2\leq x\int_0^x(v')^2\leq\|v\|_E^2\).
Integrating yields (H17). The integrals in (H18) are explicit polynomial
antiderivatives, so the endpoint and second-derivative assertions follow
by coefficient calculation. Polynomial integration by parts, with the
zero boundary values, proves the first equation in (H19). Taking \(v=Gp\),
\[
\|Gp\|_E^2=\int p\,Gp
 \leq\|p\|_{L^2}\|Gp\|_{L^2}
 \leq\|p\|_{L^2}\|Gp\|_E,
\]
which proves the norm estimate.

Cauchy--Schwarz for polynomial integrals follows by integrating the square
of \(sp+tq\) and using nonnegativity for every rational \(s,t\); the resulting
quadratic has the required discriminant inequality. Nonnegativity of the
integral of a polynomial known nonnegative on an interval follows from
finite Riemann sums and its explicit polynomial Lipschitz bound. Thus every
analytic-looking step has a finite rational-sum proof with a modulus, as in
the existing integration packet \cite{repoFramework}. The integration and
derivative identities themselves are exact finite coefficient identities.
\end{proof}

An energy name therefore defines a uniform continuous-function name:
(H17) bounds its uniform error by its energy error. The derivative need only
be an \(L^2\) name; a pointwise derivative is not inferred.

\begin{theorem}[A weak elliptic Dirichlet equation]
\label{thm:pde}
Suppose \(a\) is given by primitive-recursive rational polynomials \(a_k\)
with uniform error at most \(2^{-k}\) on \([0,1]\), and rational bounds
\[
                     0<m\leq a(x)\leq M .
\]
Suppose \(f\in L^2_P\) is given by a primitive-recursive polynomial name
\(f_k\), with \(\|f_k-f\|_{L^2}\leq2^{-k}\) and \(\|f\|_{L^2}\leq B\).
Then a primitive-recursive regular sequence of rational zero-boundary
polynomials names a unique \(u\in H_E\) such that
\[
                \int_0^1 a\,u'v'=\int_0^1f v
                         \qquad(v\in H_E).
\tag{H20}
\]
Its energy error at index \(n\) is at most \(2^{-n-2}\); the same bound
holds in uniform norm. Equation (H20) is the weak meaning of
\(-\left(a u'\right)'=f\), \(u(0)=u(1)=0\).
\end{theorem}
\begin{proof}
For \(v\in P_0\), define \(Av\) by its derivative
\[
 (Av)'=a v'-\int_0^1a(t)v'(t)\,dt,\qquad (Av)(0)=0.
\tag{H21}
\]
The derivative has mean zero, so also \((Av)(1)=0\).
Its energy interpretation is obtained by replacing \(a\) with a rational
polynomial approximation, integrating, and subtracting the mean.
The map \(g\mapsto g-\int_0^1g\) is contractive in \(L^2\), because
\[
       \int_0^1(g-\overline g)^2
                      =\int_0^1g^2-\overline g^{\,2}\leq\int_0^1g^2 .
\tag{H22}
\]
An approximation \(a_k\) consequently changes (H21) by at most
\(2^{-k}\|v\|_E\).
The rational \(U(v)=1+\|v\|_E^2\) bounds \(\|v\|_E\).
For accuracy \(2^{-r}\), take
\(k=r+\max(1,\lceil U(v)\rceil)\).
Then \(2^{-k}U(v)\leq2^{-r}\).
This constructs a primitive-recursive strong evaluator of \(A\).

For \(v,w\in P_0\), the mean in (H21) contributes nothing to
\(\int(Av)'w'\), because \(\int_0^1w'=0\). Hence
\[
 \langle Av,w\rangle_E=\int_0^1a v'w',\qquad
 \langle Av,v\rangle_E\geq m\|v\|_E^2,\qquad
 \|Av\|_E\leq M\|v\|_E.
\tag{H23}
\]
The last inequality follows from (H22) and
\(\|a v'\|_{L^2}\leq M\|v'\|_{L^2}\).
The explicit polynomial errors prove the same inequalities in the
completion. All strong-operator data have now been constructed.

Set \(b_\infty=Gf\) using (H18)'s extension. Equation (H19) gives the
strong evaluator \(Gf_k\), the required error, and \(\|b_\infty\|_E\leq B\).
Theorem~\ref{thm:hilbert-solver} produces \(Au=Gf\).
Pairing with \(v\) and using (H19), (H23) proves (H20), first for polynomials
and then for regular energy names by the displayed boundedness inequalities.
The same approximation argument defines the products and integrals in
(H20), with errors controlled by Cauchy--Schwarz.
Uniqueness and the uniform error follow from coercivity and (H17).
\end{proof}

\begin{example}[An exactly checkable continuous example]
\label{ex:pde}
Take \(a(x)=1+x\), \(f(x)=1\), \(m=1\), \(M=2\), \(B=1\).
Then \(\lambda=1/4\), \(\delta=1/8\), \(h=8\), \(R=2\), \(C=4\), and
\(N_n=8(n+5)\). Both \(A\) and \(Gf=x(1-x)/2\) have exact rational-polynomial
evaluators. Starting at \(0\), (H8) computes a polynomial of degree at most
\(N_n+1\). The classical expression for the same solution is
\[
                       u(x)=\frac{\log(1+x)}{\log 2}-x .
\]
Indeed \((1+x)u'(x)=1/\log2-(1+x)\), its derivative is \(-1\), and both
endpoints vanish. This expression is explanatory only: the polynomial
producer and its error proof do not depend on evaluating a logarithm.
The companion program checks rational energy residuals of the produced
polynomials and uses (H15) to certify them.
\end{example}

\subsection{Why weak functional data are insufficient}
\label{subsec:riesz}

The word \emph{strong} in Definition~\ref{def:strong-operator} matters.
A bilinear form given only by values on finite rational vectors need not
give a strong operator evaluator; similarly a bounded functional need not
give a norm-convergent name for its representing vector.

\begin{theorem}[Exact missing datum for a rational \(\ell^2\) vector]
\label{thm:tail-equivalence}
Let \(a_i\in\Q\) be a fixed primitive-recursive sequence whose finite squared
sums have a supplied common upper bound. The following data are equivalent:
\begin{enumerate}
\item a primitive-recursive strong \(\ell^2\) name with coordinates \(a_i\);
\item a primitive-recursive tail modulus \(\tau\) such that, for all \(k\)
and finite \(L\geq\tau(k)\),
\[
                    \sum_{\tau(k)\leq i<L}a_i^2\leq4^{-k}.
\tag{H24}
\]
\end{enumerate}
Both directions are explicit transformations without unbounded search.
\end{theorem}
\begin{proof}
Replace \(\tau(k)\) by \(\max_{j\leq k}\tau(j)\) if necessary.
The vector of the first \(\tau(k)\) coordinates is a regular name:
the squared distance between its \(k\)-th and \(\ell\)-th truncations,
\(\ell\geq k\), is bounded by (H24).
There is no need to assume \(\tau(k)\to\infty\). If a fixed coordinate is
still omitted, (H24), applied to the finite tail ending at that coordinate,
bounds its absolute value by \(2^{-k}\). Thus the represented coordinate
is \(a_i\) in either case.

Conversely, let \(p_k\) be a finite-support rational vector at distance
at most \(2^{-k}\) from the named vector. Choose \(\tau(k)\) larger than
every index in its support. On every finite coordinate interval starting at
\(\tau(k)\), \(p_k\) vanishes. Orthogonal projection onto that interval is
contractive, so the squared sum of the corresponding \(a_i\)'s is bounded
by \(4^{-k}\). The projection inequality says merely that a partial sum of
nonnegative rational squares is at most the full sum. For named coordinates
use finer approximants and rational separation. The support bound is
computed from the code \(p_k\), and hence is primitive recursive.
\end{proof}

\begin{proposition}[A bounded functional without a computable strong representative]
\label{prop:riesz-obstruction}
There is a rational primitive-recursive functional on finite rational
sequences, with norm bound \(1\), whose continuous extension to
\(\ell^2\) has no computable strong Riesz representative. In particular,
the missing tail modulus cannot be uniformly manufactured in \(\ipra\).
\end{proposition}
\begin{proof}
Fix a standard effective enumeration of machines and a primitive-recursive
pairing \(\langle e,t\rangle\). Whether machine \(e\) first halts at step
\(t\) is decided by a bounded simulation. Define
\[
 a_{\langle e,t\rangle}=
 \begin{cases}
 2^{-e-1},&\text{if \(e\) first halts at step \(t\)},\\
 0,&\text{otherwise.}
 \end{cases}
\tag{H25}
\]
This is a primitive-recursive rational sequence. Each \(e\) contributes
at most one nonzero square, so every finite squared sum is at most
\(\sum_{e\geq0}4^{-e-1}=1/3<1\). Only the finite geometric-sum inequality
is needed to prove this bound.

For finite rational \(v\), put \(L(v)=\sum_i a_i v_i\).
This is a primitive-recursive exact finite sum, and Cauchy--Schwarz gives
\(|L(v)|\leq\|v\|\). It extends to a computable functional on strong
\(\ell^2\) names: apply it to an input approximant and retain the same error.
A Riesz representative must have coordinate \(L(e_i)=a_i\).

A computable strong name for that representative would, by the proof of
Theorem~\ref{thm:tail-equivalence}, give a computable tail modulus.
For a machine \(e\), put \(N=\tau(e+2)\). A nonzero coordinate
\(\langle e,t\rangle\geq N\) would contribute
\(4^{-e-1}>4^{-e-2}\) to a finite tail, contradicting (H24).
Thus \(e\) halts exactly when one of the finitely many coordinates \(i<N\)
encoding \(e\) is nonzero. This decides the halting problem, a contradiction
by its elementary diagonal undecidability proof.
\end{proof}

This is a metatheoretic noncomputability result, not a formal independence
proof of a particular sentence of \(\ipra\).
It already blocks an unrestricted computable solution operator: take the
bilinear form \(\langle u,v\rangle\) and the functional \(L\). The weak
Lax--Milgram conclusion would require exactly the missing strong
representative. The positive PDE theorem avoids this problem by constructing
its Green representative and strong operator.

\begin{corollary}[Coercive bilinear values need not supply a strong operator]
\label{cor:bilinear-obstruction}
Even an exactly primitive-recursive bilinear form on finite rational vectors,
with coercivity bound \(1\) and boundedness constant \(2\), need not have
a computable strong representing operator.
\end{corollary}
\begin{proof}
Use the functional \(L\) just constructed and set
\(\mathfrak a(v,w)=\langle v,w\rangle+L(v)L(w)\).
This form is computed by rational finite sums. Its diagonal is at least
\(\|v\|^2\), and Cauchy--Schwarz and \(\|L\|\leq1\) give
\(|\mathfrak a(v,w)|\leq2\|v\|\|w\|\).
In the ordinary completion its representing operator is \(A=I+a\otimes a\),
where \(a\) is the vector of coordinates (H25).
Choose the index \(j\) of a known immediately halting machine and its known
halting time. Then \(a_j>0\) is a specified rational. A strong evaluator of
\(A\), applied to the finite basis vector \(e_j\), would give the computable
strong name
\[
                              a=(Ae_j-e_j)/a_j.
\]
This contradicts Proposition~\ref{prop:riesz-obstruction}.
\end{proof}

\section{Complex analysis with a radius gap and simultaneous derivative control}
\label{sec:complex}

The function-space construction in Section~\ref{sec:function-topology}
requires an explicit common modulus. For analytic families that information
can be produced from a coefficient bound and a radius gap. We give a
direct construction that controls every derivative up to a requested order,
using only rational polynomial calculations and geometric bounds.
This is a quantitative normal-family statement in a \(C^J\) metric, without
asserting the existence of a computable subsequence.

The information supplied with an analytic function is significant.
Representing a function in a continuous-function space and additionally
knowing that it is analytic is not generally equivalent to having an analytic
representation with effective radius and bound data
\cite{PaulySteinberg2018}. We therefore specify this data as part of the
presentation and show exactly what follows from it.

\subsection{A coefficient presentation}

Fix positive rational \(R\), nonnegative rational \(M\), and coefficients
\(a_n\in\C\) with a joint primitive-recursive Gaussian-rational evaluator.
The normalized coefficients
\[
                         b_n=R^n a_n
\]
have Gaussian-rational approximants to every dyadic precision, and the input
certificate asserts
\[
                             |b_n|\leq M\quad(n\in\N).
\tag{A1}
\]
This is a schema in one fixed joint evaluator. It does not enumerate all
analytic functions. The case \(M=0\) yields the zero function and can be
handled separately.

Put \(P_N(z)=\sum_{n=0}^Na_nz^n\).
For \(j\geq0\) write
\((n)_j=n(n-1)\cdots(n-j+1)\), with \((n)_0=1\) and
\((n)_j=0\) when \(n<j\). Then
\[
                 P_N^{(j)}(z)=
                   \sum_{n=j}^N(n)_j a_n z^{n-j}
\]
is a finite identity. The superscript is initially only formal polynomial
differentiation; the theorem below proves agreement with analytic
differentiation.

\begin{lemma}[A finite differentiated geometric bound]
\label{lem:derivative-geometric}
For rational \(0\leq t<1\), all \(j,L\in\N\),
\[
            \sum_{n=j}^L(n)_j t^{n-j}
                      \leq \frac{j!}{(1-t)^{j+1}}.
\tag{A2}
\]
The assertion is a finite rational inequality, with no differentiated
infinite series in its proof.
\end{lemma}
\begin{proof}
If \(L<j\), the left side is zero. Otherwise the following finite expansion
has nonnegative degrees throughout.
There are \(\binom{m+j}{j}\) ordered \((j+1)\)-tuples of nonnegative
integers with sum \(m\). This identity follows by finite induction, or by
placing \(j\) separators among \(m+j\) positions.
Expand
\((1+t+\cdots+t^L)^{j+1}\).
All its coefficients are nonnegative, and its terms of total degree
at most \(L-j\) include every such tuple. Consequently
\[
\sum_{m=0}^{L-j}\binom{m+j}{j}t^m
 \leq(1+t+\cdots+t^L)^{j+1}
 \leq(1-t)^{-j-1}.
\]
Multiply by \(j!\) and use
\((m+j)_j=j!\binom{m+j}{j}\).
If \(L<j\), the left side of (A2) is zero.
The tuple count, expansion, and geometric inequality are all finite
primitive-recursive calculations.
\end{proof}

\subsection{A uniform tail for all requested derivatives}

Fix rational \(0\leq r<R\) and an order \(J\in\N\). Define
\[
s=\frac{r+R}{2},\qquad
\rho=\frac rR,\qquad \theta=\frac rs,\qquad \sigma=\frac sR ,
\]
so \(0\leq\theta<1\), \(0<\sigma<1\). For \(0\leq j\leq J\) put
\[
 K_j=\frac{M j!}
 {s^j(1-\theta)^{j+1}(1-\sigma)},\qquad
 D_j=\frac{j!}{R^j(1-\rho)^{j+1}}.
\tag{A3}
\]
These constants are rational and explicitly computable, including at \(r=0\).

\begin{theorem}[Derivative-tail estimate and holomorphic interpretation]
\label{thm:analytic-tail}
For every \(L\geq N\geq J\), \(0\leq j\leq J\), and \(|z|\leq r\),
\[
          |P_L^{(j)}(z)-P_N^{(j)}(z)|
                          \leq K_j\sigma^{N+1}.
\tag{A4}
\]
The series thus supplies jointly primitive-recursive uniform names of
\(f,f',\ldots,f^{(J)}\) on the smaller disk, and \(f\) is holomorphic on
\(|z|<R\).
\end{theorem}
\begin{proof}
For \(n>N\geq j\), the coefficient bound gives
\[
 |(n)_ja_nz^{n-j}|
 \leq M s^{-j}(n)_j\theta^{n-j}\sigma^n.
\]
By (A2), each nonnegative factor
\((n)_j\theta^{n-j}\) is at most
\(j!/(1-\theta)^{j+1}\).
The finite sum of \(\sigma^n\), \(N<n\leq L\), is at most
\(\sigma^{N+1}/(1-\sigma)\).
Multiplication proves (A4). This also constructs the limit names:
choose \(N\) by the explicit clock below and evaluate the finitely many
coefficients to a precision whose finite weighted error is small enough.
No unbounded search for a convergent tail is involved.

To justify analytic differentiation rather than merely name a formal
derivative, first work with \(j=0,1,2\) and an intermediate disk.
For a finite polynomial, binomial expansion gives
\[
P_N(z+h)-P_N(z)-hP_N'(z)
 =h^2\sum_{n=2}^N a_n
               \sum_{\ell=2}^n\binom n\ell z^{n-\ell}h^{\ell-2}.
\]
If \(|z|\leq r\) and \(|h|\leq s-r\), the absolute value of the right side
is bounded by
\[
\frac{|h|^2}{2}\sum_{n=2}^N n(n-1)|a_n|s^{n-2}
 \leq \frac{M}{R^2(1-s/R)^3}|h|^2 .
\tag{A5}
\]
The first inequality follows termwise from
\(\binom n\ell\leq n(n-1)\binom{n-2}{\ell-2}/2\)
and \(|z|+|h|\leq s\); the second is (A2).
Pass to arbitrarily accurate polynomial approximants in (A5), with the
explicit errors from (A4). The resulting difference-quotient error is at
most a known constant times \(|h|\), proving that the named first
derivative is the complex derivative. Division by \(h\) is used only for
increments apart from zero; the remainder inequality itself holds without
that restriction.
The same argument applied to the formally differentiated polynomials,
using (A2) with the higher orders required, proves every further
derivative. For derivative order \(j\), an explicit remainder constant is
\[
             \frac{M(j+2)!}{2R^{j+2}(1-s/R)^{j+3}}.
\]
For a chosen interior point one chooses a rational smaller disk
and collar before making this argument. Thus holomorphicity and all
differentiation claims follow from finite identities and specified moduli.
\end{proof}

\paragraph{Relation with Cauchy integral presentations.}
One can instead supply a function on a certified neighborhood of
\(\{|z|\leq R\}\), a bound \(M\), a primitive-recursive contour evaluator,
and a proved Cauchy formula. The coefficient formula then gives (A1).
If the bound is available only on smaller circles, include the corresponding
radius loss. The repository's Cauchy and residue packet develops the
required finite contour estimates under explicit differentiability and
geometric hypotheses \cite{repoCauchy}; it is a proposed source of this
adapter, not a silently completed new Lean dependency.
The present theorem itself starts at the coefficient certificate and needs
no contour theorem as an unproved premise.

\subsection{A completely explicit finite \texorpdfstring{\(C^J\)}{C-J} approximation family}

Let
\[
\begin{split}
K&=\max_{j\leq J}K_j,\qquad D=\max_{j\leq J}D_j,\\
U&=\max(1,\lceil K\rceil),\qquad
V=\max(1,\lceil8D\rceil),\qquad L_0=\max(1,\lceil M\rceil),\\
h&=\ceil{1/(1-\sigma)},\\
N&=J+h(k+U+2),\qquad p=k+V,\qquad \Delta=2^{-p}.
\end{split}
\tag{A6}
\]
Use the real rational grid
\[
       G_p=\{j\Delta:j\in\Z,\ -L_0\,2^p\leq j\leq L_0\,2^p\}.
\]
For every \(N+1\) tuple \(q_0,\ldots,q_N\in G_p+iG_p\), form
\[
                         Q_q(z)=\sum_{n=0}^N q_n R^{-n}z^n .
\tag{A7}
\]
The list contains exactly
\[
                      H=(2L_0\,2^p+1)^{\,2(N+1)}
\tag{A8}
\]
polynomials. Listing all of them can be extremely expensive, but is
primitive recursive. A chosen approximant can be computed without
enumerating the entire list.

\begin{theorem}[Simultaneous derivative entropy bound]
\label{thm:analytic-entropy}
For every family member satisfying (A1), finite observations select a
polynomial \(Q_q\) in (A7) for which
\[
 \max_{0\leq j\leq J}\ 
        \sup_{|z|\leq r}|f^{(j)}(z)-Q_q^{(j)}(z)|
                        \leq2^{-k-1}.
\tag{A9}
\]
The selecting algorithm, degree, coefficient precision and cardinality
bound are primitive recursive in the supplied rational data, \(J,k\),
and the fixed coefficient evaluator.
\end{theorem}
\begin{proof}
The proof of Lemma~\ref{lem:clock}, with \(q=\sigma\), gives
\(\sigma^h\leq1/2\). Since \(K\leq U\leq2^U\),
\[
K\sigma^{N+1}\leq U\,2^{-(k+U+2)}\leq2^{-k-2}.
\tag{A10}
\]
This is a common bound for the derivative tails.

Approximate each real and imaginary part of \(b_n\), \(0\leq n\leq N\),
within \(\Delta/4\), round to a nearest grid point, and clip to
\([-L_0,L_0]\). The true component lies in this interval, so clipping cannot
increase its error. Each coordinate error is at most \(3\Delta/4<\Delta\),
and hence \(|b_n-q_n|\leq2\Delta\).
At derivative order \(j\), (A2) gives the uniform polynomial error
\[
\left|\sum_{n=j}^N(n)_j(a_n-q_nR^{-n})z^{n-j}\right|
 \leq 2\Delta R^{-j}\sum_{n=j}^N(n)_j\rho^{n-j}
 \leq2\Delta D.
\]
By \(V\geq8D\) and \(2^V\geq V\), this is at most \(2^{-k-2}\).
Add (A10) to obtain (A9).
All selection decisions are comparisons of rational observations; none
is a decision about equality or sign of an exact complex coefficient.
\end{proof}

The approximating polynomials lie in the ambient \(C^J\) space; they need not
belong to the original analytic family or obey its original coefficient
bound. Selecting representatives inside an arbitrary family would be
an additional existence problem. The theorem gives an explicit total
boundedness result for the analytic family viewed in the smaller-disk
\(C^J\) metric. It does not choose a convergent subsequence from an arbitrary
sequence in that family.

\subsection{An example with exact rational derivative checks}

Take \(R=2\), \(M=1\) and \(a_n=2^{-n}\). Then
\[
                      f(z)=\frac1{1-z/2},\qquad
          f^{(j)}(z)=\frac{j!\,2^{-j}}{(1-z/2)^{j+1}}.
\]
For rational complex \(z\) with \(|z|\leq r<2\), both the finite polynomial
jets and the right-hand expressions are exactly Gaussian rational.
The reference program compares them using squared norms and verifies
(A4), (A10) and (A9)'s error budgets at several radii and derivative orders.
The integer schedules in (A6) are intentionally conservative; their purpose
is an unconditional primitive-recursive bound, not good numerical complexity.

\subsection{The radius gap and the analytic structure cannot be discarded}

On the closed unit disk the polynomials \(z^n/n\) tend uniformly to zero,
while their derivatives \(z^{n-1}\) have uniform norm \(1\).
Thus differentiation is not continuous for the same-disk uniform metric.
The supplied smaller radius changes the conclusion: the powers of
\(\sigma<1\) absorb all polynomial derivative factors. This explains the
appearance of a collar in analytic arguments, rather than treating it as
a technical artifact.

The continuous partitions of unity from the topology section must also be
used with care. They glue approximately compatible continuous data, but
their weights are generally not holomorphic. An approximately glued result
is therefore not automatically holomorphic. If the local holomorphic
functions agree exactly on the relevant overlaps, the glued function equals
them locally and is holomorphic; otherwise a separate analytic correction
argument is needed. This is one reason to retain domain-specific
interpretation lemmas even in a common framework.

\subsection{Object statements and the Lean path}

Encode a Gaussian rational as a pair of rationals, a polynomial as a finite
coefficient list, and a jet as the list of its formal derivatives through
order \(J\). Horner evaluation, finite differentiation, factorials, powers,
grid rounding, and (A6) are fixed primitive-recursive programs.
The core object statements quantify over finite \(L,N,J\), a rational query
\(z\), and rational coefficient observations, with explicit error terms for
those observations. The majorant conclusion is a rational squared-norm
inequality; all denominators \(R,s,1-\rho,1-\theta,1-\sigma\) have proved positive
lower bounds.

The needed derivations are: the finite geometric identity; (A2) by a bounded
coefficient count; a triangle estimate for finite complex sums; the monomial
derivative identity; the split-radius majorant (A4); the dyadic clock;
the rounded-coefficient estimate; and the finite difference identity (A5).
Their induction invariants are bounds for finite partial sums or coefficient
lists, not holomorphicity assertions. After these are obtained, regularity,
extensionality, and analytic interpretation are assembled as in
Theorem~\ref{thm:compilation}. The actual object proof terms for these new
complex-analytic statements remain explicit formalization targets.

\section{Checked Lean nucleus and a concrete formalization route}
\label{sec:formalization}

The archive separates an executable, kernel-checked development from the
larger object-language implementation that remains to be written. This
section records that boundary and gives the definitions and induction
targets needed to cross it. The detailed module plan is also supplied in
\path{lean/OBJECT_LANGUAGE_PLAN.md}.

\subsection{Exactly what the delivered Lean development proves}

\path{lean/QuantitativeKernel.lean} contains 19 theorem declarations with
complete proofs. It and \path{lean/Examples.lean} were checked using Lean
4.32.0, compiler commit
\texttt{8c9756b28d64dab099da31a4c09229a9e6a2ef35}.
There are no \texttt{sorry} placeholders and no new axiom declarations.
The complete replay output is in \path{lean/verification.log}.

\begin{center}
\small
\begin{tabular}{p{0.27\textwidth}p{0.63\textwidth}}
\toprule
Delivered component & Exact status and scope\\
\midrule
\path{positive_transfer} & Checked host Lean theorem: preservation of positive
formulas under an atomic transfer contract and explicit forward/backward
witness selectors. No axiom dependencies.\\[2pt]
\path{allFin_true}, \path{anyFin_true}, \path{eval_correct} & Checked Boolean
folds and finite-domain evaluator correctness. The reported dependencies
are \texttt{propext}, and additionally \texttt{Quot.sound} for the evaluator.\\[2pt]
Rational error and clock lemmas & Checked perturbation invariant, block
halving, amplitude absorption, and the full explicit
\path{solver_budget} schedule. The dependency audit reports
\texttt{propext}, \texttt{Classical.choice}, and \texttt{Quot.sound}.\\[2pt]
New main analytic results & Human proofs with explicit arithmetic targets;
no new main theorem in this archive has been delivered as a literal
\texttt{Prf PRAAx} proof value.\\
\bottomrule
\end{tabular}
\end{center}

The arithmetic axioms in the third row are inherited from the standard
\texttt{Rat} library: even its elementary power and order lemmas report
those dependencies. Thus the arithmetic host development is not advertised
as choice-free. Its proposed \(\ipra\) implementation uses the
denominator-cleared natural inequalities below. Host metatheoretic axioms
and the axioms available to an object derivation are different questions.

The checked project imports only \texttt{Std} and
\path{Init.Data.Rat.Lemmas}; it has no Mathlib or ProveIt import.
All project sources and their transitive project dependencies are included.
The external prerequisite is the pinned Lean distribution with its standard
library. \path{lean/PROVENANCE.json} records the compiler, distribution
checksum, and inspected repository sources. The command
\texttt{bash replay.sh} in the \texttt{lean} directory builds the library,
prints the dependency audit, and checks the examples. The compiler itself
and disposable build files are not included in the archive.

\subsection{Definitions at the semantic and arithmetic boundaries}

The checked logical syntax is completely explicit.
\texttt{Env X} means \texttt{Nat -> X}; \texttt{push} places a newly bound
variable at index zero and shifts the previous environment. The relation
\texttt{Related R} compares environments pointwise.
\texttt{PosFormula Atom} has constructors for an atom, truth, falsity,
conjunction, disjunction, and the two first-order quantifiers.
\texttt{Holds} recursively interprets those constructors.
\texttt{positive\_transfer} takes maps \(X\to Y\) and \(Y\to X\)
choosing related points. Existential transfer uses the first map;
universal transfer uses the second. Different margins can be incorporated
into the source and target atomic interpretations. Neither connective nor
quantifier adds a further margin. The Boolean evaluator uses
\texttt{allFin} and \texttt{anyFin}, proved correct by induction on the
finite domain size.

These definitions are ordinary Lean types. For an arithmetic instance, a
host function parameter is replaced by a specified PR term or a finite
input table with its proved contract. The additional data for the analytic
modules are as follows; all displayed predicates have been defined
mathematically in the preceding sections.

\begin{description}[style=nextline,leftmargin=1em]
\item[Rational and dense codes.]
A rational code is a signed integer numerator and a positive natural
denominator. Gaussian rational codes are pairs of these. A metric basis
has a decidable domain of natural codes, a fixed rational distance program,
and derived metric laws. For a pre-Hilbert basis use rational vector
operations and a rational inner product. The finite radius predicate is
\[
 \operatorname{Ball}(v,w,r)\;:\Longleftrightarrow\;
 r\geq0\ \land\ \langle v-w,v-w\rangle\leq r^2,
 \tag{F1}\label{eq:formal-ball}
\]
with validity of the two vector codes included in its domain guard.
It does not require an exact rational value of a generally irrational norm.

\item[Names, moduli, and observation contracts.]
For a fixed producer \(p(a,n)\), \(\Reg_p(a)\) is the formula from
Definition~\ref{def:regular}, with universal precision variables and
explicit metric tests. \(\Eq_{p,q}(a,b)\) is the stated finer-index equality
test. A dense evaluator stores its fixed evaluation term, monotone demand
term, and joint continuity derivations. Net data comprise a positive size,
an indexed list of valid centers, and the covering formula
\eqref{eq:top-net}. An evaluator for a function family includes exactly the
three rational contracts \eqref{eq:top-name}--\eqref{eq:top-space}.

\item[Polynomial and formula codes.]
A polynomial is a finite coefficient list; multiplication is convolution,
differentiation sends coefficient \(p_{j+1}\) to \((j+1)p_{j+1}\), and
integration from zero sends \(p_j\) to coefficient \(p_j/(j+1)\) at
index \(j+1\). A fixed multisorted signature records arities, input/output
sorts, symbol evaluators, and moduli. Terms and continuous formulas are
finite trees with the constructors in Section~\ref{sec:cm-model}; finite
environments respect the recorded sorts. Finite restriction replaces each
quantifier by a minimum or maximum scan, preserving its scope.
\end{description}

In the object implementation, tuple coding collects numerical parameters
into the arguments of fixed \texttt{PTm} terms. Inequality predicates are
quoted as \texttt{Fm PTm} formulas, with positive denominators cleared.
The intended output is \(\Gamma\vdash[\mathrm{PRAAx}]\varphi\).
The context \(\Gamma\) retains every input law that has not already been
proved. Merely filling a semantic record with assumptions does not
discharge those hypotheses.

\subsection{Two fully specified quantifier-free induction targets}

The new scalar estimates already have explicit arithmetic invariants.
Write the contraction factor as \(q=a/b\), initial majorant as \(R=r/s\),
per-step error as \(p/t\), and reserve as \(e/f\), with natural numerators,
positive denominators, and \(a\leq b\). The stationary-noise condition is
\[
                    pbf\leq(b-a)et.
\]
Define a PR numerator sequence
\[
 U_0=rt,\qquad U_{j+1}=aU_j+psb^{j+1}.
\]
Then \(U_j/(stb^j)\) is the exact majorant recurrence. The induction formula
is the quantifier-free inequality
\[
              fU_j\leq rtf a^j+est b^j.
 \tag{F2}\label{eq:formal-noise-invariant}
\]
Its base case is immediate. The next-step slack is identically
\[
 a\bigl[rtf a^j+estb^j-fU_j\bigr]
       +sb^j\bigl[(b-a)et-pbf\bigr].
\]
Both bracketed slacks are nonnegative. The repository's polynomial
inequality machinery can therefore prove the step with those two explicit
weights; \texttt{PRA.qfInd} yields the invariant. A trace whose error
certificates satisfy the local recurrence is compared with this majorant
by another bounded induction. No induction on a Cauchy predicate is used.

For the clock, put \(\delta=d/b\), \(0<d\leq b\). The corresponding
natural-number invariant is
\[
             (b+jd)(b-d)^j\leq b^{j+1}.
 \tag{F3}\label{eq:formal-clock-invariant}
\]
The exact step identity is
\[
 (b-d)\bigl(b+(j+1)d\bigr)+(j+1)d^2=b(b+jd).
\]
Multiply by \((b-d)^j\), discard the nonnegative term, and apply the
previous invariant. If desired, use a separate nonnegative parameter
\(a\) with \(a+d=b\) to eliminate truncated subtraction from the identity.
When \(hd\geq b\), specialization and positive cancellation give
\(2(b-d)^h\leq b^h\), the halving step. The elementary induction
\(C\leq2^C\) supplies the remaining precision offsets.

The checked declaration \texttt{solver\_budget} verifies the resulting
rational statement:
\[
\begin{gathered}
 0\leq\delta\leq1,\quad h\delta\geq1,\quad2\lambda\leq\delta C,\\
 x_0\leq R,\qquad
 x_{j+1}\leq(1-\delta)x_j+2\lambda2^{-(n+3+C)}
 \quad\Longrightarrow\quad
 x_{h(n+R+3)}\leq2^{-n-2}.
\end{gathered}
\]
Here \(h,C,R,n\) are natural numbers and \(x_j\) is a rational majorant.
The Hilbert application must establish its local radius contract
\eqref{eq:formal-ball}; this checked scalar theorem does not by itself
formalize an infinite-dimensional operator.

\subsection{Formalizing the main constructions in dependency order}

\paragraph{Completion and categorical operations.}
First quote the metric and joint evaluator laws, prove their substitution
instances, and apply the explicit demand in
Theorem~\ref{thm:cm-extension}. To derive preservation of \(\Reg\), introduce
arbitrary precision variables, instantiate the joint law, and combine its
rational inequalities. To derive preservation of \(\Eq\), use the finer
input indices and rational separation. These proofs do not apply
quantifier-free induction to the quantified formulas \(\Reg\) or \(\Eq\).
Product and change-of-basis maps use the same rules with maxima of demands.
Only after these obligations are proved should the fixed numerical family
be passed to \texttt{SetoidInterpretation.Presentation} and
\texttt{SetoidInterpretation.sound} \cite{repoSetoid}.

\paragraph{Function-space approximation and gluing.}
Define the rational positive part, the weights \(w_i\), partial sums
\(W_\ell\), and coefficient reconstruction. Prove the denominator bound
\(W_n\geq r/2\) from a bounded covering witness before invoking division.
For a fixed query \(x\) and evaluation precision \(p\), the central
sum invariant is
\[
 \left|\sum_{i<\ell}w_i(x)
            \bigl(h_i-E(a,x,p)\bigr)\right|
       \leq\bigl(2^{-k-1}+2^{-p}\bigr)W_\ell(x).
\]
The coefficient and support cases are decidable rational comparisons.
At the final sum, division yields \eqref{eq:top-object-target}.
Mixed-radix list coding proves the cardinality bound independently.
Gluing uses the same nonnegative-sum invariants and the already supplied
local bounds. None of these steps needs a decision about which arbitrary
profile classes are inhabited.

\paragraph{Continuous formulas and finite restrictions.}
Implement rational \(\min\), \(\max\), and \(1-u\), then finite extrema
with a first attaining index. The scan invariant says that the current
value is attained and is a lower, respectively upper, bound for the
processed prefix; its bounded universal part is Booleanized.
Prove that two finite lists related by the stated atomic-error relation
have extrema within that error. External structural induction on the
\emph{fixed} formula then constructs the arithmetic derivation of
Theorem~\ref{thm:cm-backforth}. Net comparison gives the coherent name in
Theorem~\ref{thm:cm-values}. The checked \texttt{positive\_transfer} is the
logical core of this argument; the numeric extrema lemma is additional
work, not a result already contained in that Boolean theorem.

\paragraph{The Hilbert solver and the Dirichlet application.}
Prove the dense inner-product identities and the finite Cauchy--Schwarz
inequality. Express approximate contraction as the rational rule
\[
 \operatorname{Ball}(v,w,r)\Longrightarrow
 \operatorname{Ball}\bigl(F_n(v),F_\ell(w),
                    qr+\eta_n+\eta_\ell\bigr).
 \tag{F4}\label{eq:formal-local-contraction}
\]
Derive this rule from the finite operator contracts; auxiliary approximation
errors are removed by the explicit rational separation lemma. Its
application to the ball, two-run, and adjacent-state invariants (H12)--(H14)
uses the same quantifier-free majorant calculations as
\eqref{eq:formal-noise-invariant}. The output name and its residual are
then proved coherent before the semantic fixed point is introduced.

For the PDE, define \(\operatorname{Int}(p)=\sum_jp_j/(j+1)\),
\(P_0\) by its two rational endpoint equations, and
\(\operatorname{Energy}(p,q)=\operatorname{Int}(p'q')\).
Polynomial antiderivatives implement \(G\) exactly. The approximating
operator has derivative
\(a_kv'-\operatorname{Int}(a_kv')\) and zero initial value.
Coefficient identities prove the endpoints, mean-subtraction identity,
and Green identity. Finite Riemann estimates with the polynomial modulus
prove the required nonnegative-integral inequalities. These derive the
strong approximation, coercivity, and norm contracts needed by
Theorem~\ref{thm:hilbert-solver}. Extension of the final bilinear identity
uses its explicit error bounds; no general Riesz or weak compactness
theorem is an omitted prerequisite.

\paragraph{Complex derivatives and finite analytic covers.}
Use finite Gaussian-rational coefficient lists and the derivative operation
defined above. Prove the finite differentiated-geometric inequality
of Lemma~\ref{lem:derivative-geometric}, including the \(L<j\) zero case,
by its finite coefficient-counting identity. This supplies the uniform
tail and coefficient-perturbation bounds for all \(j\leq J\).
The constants \(K_j,D_j\) and the positive radius gaps are rational;
Booleanized maxima over \(j\leq J\) allocate one common precision.
The clock calculations reuse the same block-halving arithmetic.
Finite binomial expansion proves the quadratic difference-quotient remainder
(A5); substitution of arbitrarily fine polynomial observations identifies
the formal derivatives with the complex derivatives. Finally encode the
coefficient grid by mixed radix and quote the explicit count (A8).
This proves the finite \(C^J\) cover, without selecting a convergent
subsequence or assuming that its centers belong to the input family.

\subsection{Repository interfaces and the final acceptance criterion}

The following existing imports provide the intended object backend; they
are not dependencies of the standalone checked project:
\begin{center}
\small
\begin{tabular}{p{0.42\textwidth}p{0.48\textwidth}}
\toprule
Within \path{IPRADerivations/} & Intended use\\
\midrule
\path{Foundations/ObservationComposition} &
\path{compose_obs}: combine finite queries and their maximum demand.\\
\path{Foundations/FiniteReplay} &
\path{replay}: preserve a completed trace when the input table is extended.\\
\path{Foundations/SetoidInterpretation} &
\path{sound}: transport a fixed presented first-order theory after its
domain, congruence, and axiom obligations are proved.\\
\path{Foundations/ContractionIterates} &
Existing exact half-contraction induction; extend it with the stationary
noise invariant rather than silently treating it as already present.\\
\path{Foundations/DiscreteGronwall} &
Existing denominator-cleared recurrence proofs. Its growth-factor theorem
is not itself the new contraction-reserve theorem.\\
\path{Analysis/ErrorBudgets} &
Finite sums, products, triangle estimates, and separated reciprocals.\\
\bottomrule
\end{tabular}
\end{center}

A completed integration should export explicit \texttt{Prf PRAAx}
declarations for the finite approximation, coherence, and residual targets,
retain their actual hypotheses, and audit their transitive proof
dependencies. Semantic adequacy should then identify each endpoint with
the represented theorem. A successful host Lean build is valuable evidence
for the arithmetic and structural arguments; it becomes evidence of
object-\(\ipra\) provability only when those additional proof values have
actually been constructed and checked.

\section{What has been resolved, what remains, and how to reuse the work}
\label{sec:synthesis}

\subsection{A common framework with substantive domain theorems}

The unifying object is a \emph{specified presentation}, not a discrete
replacement of the mathematical object. A name may represent a point in an
infinite-dimensional Hilbert space, a function on an arbitrary presented
compact metric space, or a germ of a holomorphic function. What is finite
is a requested observation and the computation proving its error bound.
Four recurring ingredients determine the successful cases:
\begin{enumerate}
\item A dense family of finite objects with primitive-recursive operations.
\item A modulus converting the desired output precision to input demands.
\item A stability inequality converting approximate local work into a
coherent global result.
\item A semantic theorem identifying the represented result with the
intended continuous object.
\end{enumerate}
None of the four is automatic. The report produces them for concrete
classes rather than adding them as names for the desired conclusion.

\begin{longtable}{p{.25\textwidth}p{.36\textwidth}p{.30\textwidth}}
\toprule
Difficulty & Result supplied here & Exact remaining boundary\\
\midrule
\endfirsthead
\toprule
Difficulty & Result supplied here & Exact remaining boundary\\
\midrule
\endhead
An operation is given only on dense finite data &
Theorem~\ref{thm:cm-extension} constructs its extension and proves independence
of names from a joint modulus. &
The dense evaluator and its joint modulus must actually be supplied.\\[2mm]
A change of presentation appears to change the object &
Proposition~\ref{prop:cm-products-density} constructs inverse name maps from
a dense isometry and an approximate inverse selector. &
A mere assertion of density does not uniformly provide that selector.\\[2mm]
Continuous formulas contain alternating quantifiers &
Theorems~\ref{thm:cm-backforth} and \ref{thm:cm-values} construct finite
approximations to their values, retaining alternation. &
Numerical approximation does not decide exact zero truth or prove
logical compactness.\\[2mm]
A family of continuous functions needs a common finite representation &
Theorem~\ref{thm:quant-ascoli} constructs tent interpolants and an explicit
entropy bound. Proposition~\ref{prop:ascoli-converse} gives a quantitative
converse. &
A common modulus and effective domain nets are required; pointwise
continuity alone is insufficient.\\[2mm]
Local formulas must be assembled without dividing by an uncertified
denominator &
Theorem~\ref{thm:guarded-gluing} derives a denominator lower bound, an
evaluator, a Lipschitz constant, and the gluing error. &
Approximate continuous gluing need not preserve holomorphicity.\\[2mm]
Finite approximations need to define one solution &
Theorem~\ref{thm:hilbert-solver} compares two finite runs directly.
The minimax example also produces a coherence modulus from a proved gap. &
Unstable approximate solutions need not admit any effective coherent choice.\\[2mm]
An infinite-dimensional equation needs an actual solution producer &
Theorem~\ref{thm:hilbert-solver} derives the contraction gap, clock and
precision; Theorem~\ref{thm:pde} constructs a weak elliptic solution. &
The strong operator data and positive coercivity margin are essential.\\[2mm]
Weak Hilbert data lack a norm name &
Theorem~\ref{thm:tail-equivalence} identifies the exact tail information
needed; Proposition~\ref{prop:riesz-obstruction} proves it cannot be
manufactured in general. &
No blanket Riesz or weak-data Lax--Milgram algorithm follows.\\[2mm]
Differentiation consumes uncontrolled precision &
Theorems~\ref{thm:analytic-tail} and \ref{thm:analytic-entropy} construct
simultaneous derivative approximants with explicit degree and precision. &
A radius gap and coefficient bound are retained; there is no same-disk
uniform differentiation modulus.\\
\bottomrule
\end{longtable}

\subsection{The source audit changes some apparent problems}

The review found that several general proof-transport ideas are already
implemented in the current source, even where old packet README files still
describe uncompiled proposals. Conversely, not every declaration located
under \sourcepath{Derivations/} has an object-theory conclusion.
The inspected \sourcepath{Foundations/FiniteUseSignObstruction.lean}
expressly concerns external Lean propositions about oracle computation.
\sourcepath{Foundations/CoherenceObstruction.lean} supplies finite
arithmetic diagonal derivations while explicitly leaving its full
simulator-index noncomputability argument outside the formalized result.
The source ledger records these distinctions. Merely counting files,
theorem declarations, or compilation successes would obscure the real
verification boundary.

The work therefore gives priority to new mathematical instances and to
closing concrete representation obligations. A new general-purpose
proof-graph wrapper would have contributed less than the explicit PDE
operator construction, the continuous function-space cover, or the
tail-modulus equivalence.

\subsection{Formal and executable deliverables}

\begin{center}
\small
\begin{tabular}{p{.29\textwidth}p{.62\textwidth}}
\toprule
Artifact & Status and purpose\\
\midrule
\sourcepath{report.tex}, included sections, and \sourcepath{report.pdf} &
Full human mathematical proofs, logical interpretation, source discussion,
and formalization plan.\\
\sourcepath{lean/QuantitativeKernel.lean} &
19 fully proved host Lean lemmas, compiled with Lean 4.32.0; no
\texttt{sorry} and no newly declared axioms.\\
\sourcepath{lean/Audit.lean} and \sourcepath{lean/verification.log} &
Actual axiom audit and replay output. Some rational lemmas inherit standard
Lean metatheoretic axioms, including \texttt{Classical.choice}.\\
\sourcepath{lean/OBJECT_LANGUAGE_PLAN.md} &
Exact denominator-cleared invariants, source interfaces and remaining
\texttt{Prf PRAAx} targets.\\
\sourcepath{code/} &
Four standard-library Python programs using exact rational arithmetic.\\
\sourcepath{evidence/} &
Source audit, finite check results, and document/build verification.\\
\bottomrule
\end{tabular}
\end{center}

The finite examples include the alternating minimax formula and its optimizer;
continuous interpolation, clipping and guarded gluing; a genuinely
infinite diagonal Hilbert equation evaluated through finite support;
the variable-coefficient elliptic equation of Example~\ref{ex:pde}; and
complex derivative jets. Their success is evidence about implementations
and constants. The universal theorems rest on the printed mathematical
proofs, and their fully internal formalizations remain the enumerated targets.

\subsection{An ordered path into the existing repository}

The most useful first addition to the object library is the
denominator-cleared perturbed-contraction and clock package.
It has a small arithmetic dependency surface and immediately supports the
Hilbert, analytic-tail and stability applications.
Next come rational polynomial energy operations and the finite
two-run comparison. The operator and PDE endpoints can then be stated in
the same named-object style as existing analysis derivations.

In parallel, the partition-of-unity rational evaluator can be built on the
existing finite-net and rational-arithmetic infrastructure. That supplies
a reusable function-space layer. A continuous-formula module should
then implement fixed syntax and finite quantifier restrictions explicitly,
instead of treating a finite net as a finite submodel. Finally, the
split-radius complex estimates provide a specialized \(C^J\) family
presentation to which both function-space and continuous-logic observations
can be applied.

Each integration step should end with an explicit closed arithmetic
sentence or open contract and a kernel-checked
\(\texttt{Prf PRAAx}\) term, followed by an axiom audit.
The semantic completion or PDE explanation belongs beside that endpoint,
so a future user can see both what was proved internally and what continuous
mathematics it represents.

\begin{thebibliography}{99}
\addcontentsline{toc}{section}{References}

\bibitem{repoMain}
ProveIt project.
\emph{Primitive recursive arithmetic: README, unification plan, and endpoint receipts}.
Revision \texttt{5f3be02388442f16487d2dc5c2529961dd783d2b}, inspected
20 November 2026.
\href{https://github.com/VladimirReshetnikov/ProveIt/tree/5f3be02388442f16487d2dc5c2529961dd783d2b/Logic/PrimitiveRecursiveArithmetic}{Pinned source directory}.
The companion source ledger records the narrower set of files actually read;
this report is not an audit of the entire repository.

\bibitem{repoSetoid}
ProveIt project.
\emph{SetoidInterpretation.lean}.
\sourcepath{Derivations/Lean/IPRADerivations/Foundations/SetoidInterpretation.lean}.
Same pinned revision as \cite{repoMain}.
\href{https://github.com/VladimirReshetnikov/ProveIt/blob/5f3be02388442f16487d2dc5c2529961dd783d2b/Logic/PrimitiveRecursiveArithmetic/Derivations/Lean/IPRADerivations/Foundations/SetoidInterpretation.lean}{Source}.

\bibitem{repoExtensional}
ProveIt project.
\emph{ExtensionalComposition.lean}.
\sourcepath{Derivations/Lean/IPRADerivations/Foundations/ExtensionalComposition.lean}.
Same pinned revision.
\href{https://github.com/VladimirReshetnikov/ProveIt/blob/5f3be02388442f16487d2dc5c2529961dd783d2b/Logic/PrimitiveRecursiveArithmetic/Derivations/Lean/IPRADerivations/Foundations/ExtensionalComposition.lean}{Source}.

\bibitem{repoUniversal}
ProveIt project.
\emph{UniversalEvaluatorObstruction.lean}.
\sourcepath{Derivations/Lean/IPRADerivations/Foundations/UniversalEvaluatorObstruction.lean}.
Same pinned revision.
\href{https://github.com/VladimirReshetnikov/ProveIt/blob/5f3be02388442f16487d2dc5c2529961dd783d2b/Logic/PrimitiveRecursiveArithmetic/Derivations/Lean/IPRADerivations/Foundations/UniversalEvaluatorObstruction.lean}{Source}.

\bibitem{repoFramework}
ProveIt research packet.
\emph{A framework for finitary mathematics in intuitionistic primitive
recursive arithmetic}.
\sourcepath{docs/proofs/foundations/ipra_framework/}.
In particular, sections 04, 06 and 06b of the packet's paper were inspected.
Same pinned revision.
\href{https://github.com/VladimirReshetnikov/ProveIt/tree/5f3be02388442f16487d2dc5c2529961dd783d2b/Logic/PrimitiveRecursiveArithmetic/docs/proofs/foundations/ipra_framework}{Packet}.

\bibitem{repoCoherent}
ProveIt research packet.
\emph{Coherent finite methods for intuitionistic primitive recursive arithmetic}.
\sourcepath{docs/proofs/foundations/ipra_coherent_methods_2026-11-15/}.
The README was inspected at the pinned revision.
\href{https://github.com/VladimirReshetnikov/ProveIt/tree/5f3be02388442f16487d2dc5c2529961dd783d2b/Logic/PrimitiveRecursiveArithmetic/docs/proofs/foundations/ipra_coherent_methods_2026-11-15}{Packet}.

\bibitem{repoCauchy}
ProveIt research packet.
\emph{Cauchy's integral theorem and the residue theorem in iPRA}.
\sourcepath{docs/proofs/analysis/cauchy_residue_ipra/}.
The README was inspected at the pinned revision.
\href{https://github.com/VladimirReshetnikov/ProveIt/tree/5f3be02388442f16487d2dc5c2529961dd783d2b/Logic/PrimitiveRecursiveArithmetic/docs/proofs/analysis/cauchy_residue_ipra}{Packet}.

\bibitem{bbhu2008}
I. Ben Yaacov, A. Berenstein, C. W. Henson, and A. Usvyatsov.
\emph{Model theory for metric structures}.
In \emph{Model Theory with Applications to Algebra and Analysis},
vol.~2, London Mathematical Society Lecture Note Series 350,
Cambridge University Press, 2008, pp.~315--427.
\href{https://doi.org/10.1017/CBO9780511735219.011}{doi:10.1017/CBO9780511735219.011}.
\href{https://math.univ-lyon1.fr/~begnac/articles/mtfms.pdf}{Author manuscript}.
Used for background on metric structures and continuous semantics;
no compactness theorem from this source is imported into \(\ipra\).

\bibitem{BrattkaYoshikawa2006}
V. Brattka and A. Yoshikawa.
\emph{Towards computability of elliptic boundary value problems in
variational formulation}.
Journal of Complexity 22(6) (2006), 858--880.
\href{https://doi.org/10.1016/j.jco.2006.04.007}{doi:10.1016/j.jco.2006.04.007}.
The bibliographic record and abstract establish the prior computable-analysis
context. The quantitative proofs in the present report are supplied in full
and do not depend on an uninspected theorem from that paper.

\bibitem{PaulySteinberg2018}
A. Pauly and F. Steinberg.
\emph{Comparing representations for function spaces in computable analysis}.
Theory of Computing Systems 62 (2018), 557--582.
\href{https://doi.org/10.1007/s00224-016-9745-6}{doi:10.1007/s00224-016-9745-6}.
\href{https://arxiv.org/abs/1512.03024}{arXiv:1512.03024}.

\bibitem{Specker1949}
E. Specker.
\emph{Nicht konstruktiv beweisbare S\"atze der Analysis}.
Journal of Symbolic Logic 14(3) (1949), 145--158.
\href{https://doi.org/10.2307/2267043}{doi:10.2307/2267043}.

\bibitem{Kreuzer2014}
A. P. Kreuzer.
\emph{From Bolzano--Weierstra{\ss} to Arzel\`a--Ascoli}.
Mathematical Logic Quarterly 60(3) (2014), 177--183.
\href{https://doi.org/10.1002/malq.201200076}{doi:10.1002/malq.201200076}.
\href{https://arxiv.org/abs/1205.5429}{arXiv:1205.5429}.
The source's reverse-mathematical equivalences concern its specified base
theories; none is asserted to be an equivalence over \(\ipra\) here.

\end{thebibliography}

\end{document}
